\documentclass[10pt]{article}
\usepackage[a4paper,margin=0.75in]{geometry}
\usepackage{booktabs,longtable,array,enumitem}
\usepackage[hidelinks]{hyperref}
\setlength{\parindent}{0pt}\setlength{\parskip}{5pt}
\renewcommand{\arraystretch}{1.12}
\newcommand{\plainbox}[1]{\par\vspace{2pt}\noindent\fbox{\parbox{\dimexpr\linewidth-2\fboxsep-2\fboxrule\relax}{#1}}\par\vspace{2pt}}
\begin{document}

{\Large\textbf{Kedarnath Yatra Workers Survey: Final Questionnaire and Variable List}}\par
{\large Version 2, worked example. Read-aloud interview, tablet-based, one visit}\par\vspace{2pt}\hrule\vspace{6pt}

\plainbox{\textbf{What this is.} The final list of questions and variables for the survey of Yatra workers. The main aim is to measure
\textbf{vulnerability to poverty} (VEP and multidimensional poverty). A short employment-quality block and a short tasks block are added.
The tasks block only gathers the facts the later paper on job transfer needs. There are \textbf{193 asked questions}, 14 recorded automatically or by the enumerator,
and 143 variables built afterwards in Stata. The companion file \texttt{variable\_dictionary.csv} lists everything in one table.}

\section*{1. Rules the questions follow}
\begin{itemize}[leftmargin=1.2em,itemsep=1pt]
\item \textbf{No rating questions.} Every answer is a fact: yes or no, a count, an amount, or a choice from a card. Nobody is asked to score anything on a scale.
\item \textbf{Aggregate household questions}, no household roster. The respondent is the worker. Consumption is for the household's usual place of living.
\item \textbf{Codes.} 98 = prefers not to say (schooling, ropeway view). 97 = don't know (job-quality questions). Blank = the question was skipped by a rule shown in the table.
\item \textbf{Recall.} 30 days for frequent spending, 12 months for infrequent spending, health, income and shocks, 15 days for recent illness.
\item Questions are read aloud in Hindi, Garhwali or Nepali. Cards are used for long answer lists.
\end{itemize}

\section*{2. Time budget}
Interview length is NOT estimated here. The module list below gives item counts only; the pretest measures real durations from the tablet's own start and end timestamps.
{\small
\begin{longtable}{llrr}
\toprule
\textbf{Module} & \textbf{Name} & \textbf{Items} & \textbf{Minutes} \\
\midrule
\endfirsthead
\toprule
\textbf{Module} & \textbf{Name} & \textbf{Items} & \textbf{Minutes} \\
\midrule
\endhead
P & Cover page, consent and paradata & 14 \\
A & Respondent and household & 20 \\
B & Work and work history & 15 \\
C & Monthly calendar of work and income & 8 \\
D & Migration, home place and remittances & 17 \\
E & Household consumption & 27 \\
F & Housing, amenities and assets & 17 \\
G & Finance, insurance and schemes & 16 \\
H & Health & 10 \\
I & Food security, shocks and coping & 10 \\
J & Ropeway & 5 \\
K & Job quality (Apablaza et al. 2026 questions) & 22 \\
L & Tasks and skills (short block) & 17 \\
M & Access to common resources (short block) & 9 \\
 & \textbf{Total} & \textbf{207} \\
\bottomrule
\end{longtable}
}

Interview length is measured from the tablet's own start and end timestamps at the pretest; no estimate is claimed here.
With 4 enumerators at 5 to 8 interviews a day, one round of about 10 days gives roughly 200 to 320 interviews.

\section*{3. Consent script (read before Module A)}
\textit{``We are doing a study on the livelihoods of people who work on the Yatra route. It takes about 40 minutes. You can stop at any time and you can skip any question. Nothing you say will be linked to your name.
Do you agree to take part?''} Record consent (variable \texttt{consent}); stop the interview if the answer is no.

\subsection*{Module P: Cover page, consent and paradata}
\begin{quote}\small\textbf{Read this out, word for word:}\\[4pt]
\textit{We are doing a study on the livelihoods of people who work on the Yatra route. I would like to ask you some questions about your work, your household and your spending.\\[6pt]Taking part is your choice. You can stop at any time, and you can skip any question you do not want to answer. Nothing you tell me will be linked to your name, and nothing you say will affect your work here or any government benefit.\\[6pt]If you say no, I will not ask you any questions. I will write down only that you said no, the date and time, and what was happening at the time, such as crowding or weather. I will not write down your name, your sex, your age, or where you are. No answers are recorded.}\end{quote}
{\scriptsize
\begin{longtable}{p{0.8cm}p{6.4cm}p{3.6cm}p{5.8cm}}
\toprule
\textbf{No.} & \textbf{Question or instruction} & \textbf{Answers} & \textbf{Variable, skip rule, source} \\
\midrule
\endfirsthead
\toprule
\textbf{No.} & \textbf{Question or instruction} & \textbf{Answers} & \textbf{Variable, skip rule, source} \\
\midrule
\endhead
P1 & Assigned by the tablet. & Automatic / enumerator & \texttt{resp\_id}\newline Source: Design \\
P2 & Set automatically from the device ID. & Automatic / enumerator & \texttt{enum\_id}\newline Source: Design \\
P3 & Set automatically from the device route setting. & Automatic / enumerator & \texttt{site}\newline Source: Project design (two-route stratification) \\
P4 & Recorded automatically. & Automatic / enumerator & \texttt{form\_build}\newline Source: Project design (build provenance) \\
P5 & Recorded automatically. & Automatic / enumerator & \texttt{interview\_date}\newline Source: Design \\
P6 & Recorded automatically. & Automatic / enumerator & \texttt{gps\_lat}\newline Source: Design \\
P7 & Recorded automatically. & Automatic / enumerator & \texttt{gps\_lon}\newline Source: Design \\
P8 & Recorded from the landing screen's consent buttons. & Automatic / enumerator & \texttt{consent}\newline Source: Design \\
P9 & Pick what was happening at the time. Do not ask the person. & Automatic / enumerator & \texttt{obs\_environment}\newline \textit{Skip: Ask only if consent = 0}\newline Source: Design: refusal record, limited to environmental conditions \\
P10 & Recorded automatically. & Automatic / enumerator & \texttt{gps\_accuracy\_m}\newline Source: Design \\
P11 & Recorded automatically. & Automatic / enumerator & \texttt{gps\_fix\_s}\newline Source: Design \\
P12 & Recorded automatically. & Automatic / enumerator & \texttt{gps\_error}\newline Source: Design \\
P13 & Recorded automatically (start to end). & Automatic / enumerator & \texttt{interview\_duration\_min}\newline Source: Design \\
P14 & Recorded automatically (module timestamps). & Automatic / enumerator & \texttt{dur\_tasks\_min}\newline Source: Design \\
\bottomrule
\end{longtable}
}
\subsection*{Module A: Respondent and household}
\begin{quote}\small\textit{Read aloud: First, a few things about you and the people who live in your household.}\end{quote}
{\scriptsize
\begin{longtable}{p{0.8cm}p{6.4cm}p{3.6cm}p{5.8cm}}
\toprule
\textbf{No.} & \textbf{Question or instruction} & \textbf{Answers} & \textbf{Variable, skip rule, source} \\
\midrule
\endfirsthead
\toprule
\textbf{No.} & \textbf{Question or instruction} & \textbf{Answers} & \textbf{Variable, skip rule, source} \\
\midrule
\endhead
A1 & How old are you? (completed years) & Number & \texttt{age}\newline Source: Standard; sensitivity covariate in the VEP model (Azeem et al. 2016 structure) \\
A2 & Record sex; ask if unsure: are you male or female? & 0 Male; 1 Female & \texttt{female}\newline Source: Standard household-survey item \\
A3 & Who is the head of your household — is it you, or someone else? If someone else, how are they related to you? & 1 Self (respondent is the head); 2 Husband; 3 Wife; 4 Father; 5 Mother; 6 Son; 7 Daughter; 8 Brother; 9 Sister; 10 Other male relative; 11 Other female relative; 12 Other, non-relative (male); 13 Other, non-relative (female) & \texttt{hoh\_relation}\newline Source: Standard household-survey item \\
A4 & What is your native language, the language you first learned at home? & 1 Garhwali; 2 Kumaoni; 3 Hindi; 4 Nepali; 5 Bhojpuri; 6 Assamese; 7 Bengali; 8 Bodo; 9 Dogri; 10 Gujarati; 11 Kannada; 12 Kashmiri; 13 Konkani; 14 Maithili; 15 Malayalam; 16 Manipuri (Meitei); 17 Marathi; 18 Odia; 19 Punjabi; 20 Sanskrit; 21 Santali; 22 Sindhi; 23 Tamil; 24 Telugu; 25 Urdu; 96 Other Indian language; 97 Other / foreign language & \texttt{native\_language}\newline Source: Project design (language proxy for regional/socioeconomic origin; avoids a sensitive caste question and covers out-of-country migrants) \\
A5 & Which language is it? Write down exactly what they say. &  & \texttt{native\_language\_other}\newline \textit{Skip: Ask only if native\_language is 96 (other Indian language) or 97 (other / foreign language)}\newline Source: Project design (the language item stands in for caste; an uncoded 'other' is a real loss) \\
A6 & Where is your permanent home? & 1 Local (same district); 2 Other Uttarakhand district; 3 Other Indian state; 4 Nepal & \texttt{origin}\newline Source: Project design (distance bucket); NOT an NSS classification — the earlier citation to one was withdrawn 2026-09-28, no NSS schedule is held by this project \\
A7 & Which state or union territory is your permanent home in? & 1 Andhra Pradesh; 2 Arunachal Pradesh; 3 Assam; 4 Bihar; 5 Chhattisgarh; 6 Goa; 7 Gujarat; 8 Haryana; 9 Himachal Pradesh; 10 Jharkhand; 11 Karnataka; 12 Kerala; 13 Madhya Pradesh; 14 Maharashtra; 15 Manipur; 16 Meghalaya; 17 Mizoram; 18 Nagaland; 19 Odisha; 20 Punjab; 21 Rajasthan; 22 Sikkim; 23 Tamil Nadu; 24 Telangana; 25 Tripura; 26 Uttar Pradesh; 28 West Bengal; 29 Andaman and Nicobar Islands; 30 Chandigarh; 31 Dadra and Nagar Haveli and Daman and Diu; 32 Delhi; 33 Jammu and Kashmir; 34 Ladakh; 35 Lakshadweep; 36 Puducherry & \texttt{home\_state}\newline Source: Standard state/UT classification; needed to apply the right poverty line per respondent \\
A8 & Is that home in a village, or in a town or city? & 1 Village (rural); 2 Town or city (urban) & \texttt{home\_rural\_urban}\newline Source: Rural and urban poverty lines (Sethu et al. 2024) \\
A9 & Are you currently married, never married, or widowed/divorced/separated? & 1 Currently married; 2 Never married; 3 Widowed/divorced/separated & \texttt{marital\_status}\newline Source: Standard household-survey item \\
A10 & Do you know exactly how many years of schooling you completed, counting from class 1? & 1 Yes, 0 No & \texttt{knows\_years\_schooling}\newline Source: Standard household-survey item \\
A11 & How many years of schooling did you complete in total? (0 if none) & Whole number & \texttt{years\_schooling}\newline \textit{Skip: Ask only if knows\_years\_schooling = 1}\newline Source: NITI Aayog National MPI (years of schooling); Chaudhuri et al. 2002 covariate \\
A12 & Which is closest: no formal education; some schooling but did not finish primary; finished primary but not secondary; or finished secondary or higher? & 1 No formal education; 2 Some schooling, did not complete primary; 3 Completed primary but not secondary; 4 Completed secondary or higher; 98 Prefer not to say & \texttt{education\_level\_cat}\newline \textit{Skip: Ask only if knows\_years\_schooling = 0}\newline Source: NITI Aayog National MPI (years of schooling); Chaudhuri et al. 2002 covariate \\
A13 & Apart from school and college, have you ever completed any training course or apprenticeship? & 1 Yes, 0 No & \texttt{training\_received}\newline Source: STEP module 2; Chaudhuri et al. 2002 adaptive-capacity covariate \\
A14 & How many members live in your household, including you? & Whole number & \texttt{hhsize}\newline Source: Chaudhuri et al. 2002; used to get per-capita consumption \\
A15 & Counting everyone in your household aged 10 or above, has at least one of them finished six years of schooling or more? & 1 Yes, 0 No & \texttt{any\_member\_6yr\_schooling}\newline \textit{Skip: Ask only if the respondent's own schooling does not already settle it: years\_schooling < 6, or the bracket is not 'completed secondary or higher'}\newline Source: NITI Aayog National MPI, Years of Schooling (weight 1/6) \\
A16 & How many of them are children under 6 years old? & Whole number & \texttt{n\_children\_u6}\newline Source: Adult-equivalent scales (Claro et al. 2010; Awuni et al. 2023); NITI Aayog National MPI nutrition age range \\
A17 & And how many are aged 6 to 14? & Whole number & \texttt{n\_children\_6\_14}\newline Source: NITI Aayog National MPI school-attendance indicator \\
A18 & How many of these children go to school now? & Whole number & \texttt{n\_children\_in\_school}\newline \textit{Skip: Ask only if n\_children\_6\_14 > 0}\newline Source: NITI Aayog National MPI school-attendance indicator \\
A19 & How many of them, including you, earn money? & Whole number & \texttt{n\_earners}\newline Source: Dependency measure used in VEP studies (Azeem et al. 2016) \\
A20 & Among those who earn money, are you the main income earner in your household? & 1 Yes, 0 No & \texttt{main\_income\_earner}\newline \textit{Skip: Ask only if n\_earners > 1 (if the respondent is the only earner, this is automatic)}\newline Source: Apablaza et al. 2026, Appendix 2 Q4 (adapted to a single-respondent report: 'who is the main contributor' becomes 'are you the main contributor') \\
\bottomrule
\end{longtable}
}
\subsection*{Module B: Work and work history}
\begin{quote}\small\textit{Read aloud: Now about the work you do here on the Yatra route, and the work you did before this.}\end{quote}
{\scriptsize
\begin{longtable}{p{0.8cm}p{6.4cm}p{3.6cm}p{5.8cm}}
\toprule
\textbf{No.} & \textbf{Question or instruction} & \textbf{Answers} & \textbf{Variable, skip rule, source} \\
\midrule
\endfirsthead
\toprule
\textbf{No.} & \textbf{Question or instruction} & \textbf{Answers} & \textbf{Variable, skip rule, source} \\
\midrule
\endhead
B1 & What is your main work during the Yatra season? Ask what they do and code the closest group; do not read the list aloud unless needed. & 1 Pony/mule worker (works the animals, does not own them); 2 Pony/mule owner (owns the animals); 3 Porter (carries on own back/shoulders - goods, luggage, or a person in a kandi); 4 Palki / dandi bearer (carries a person on a palanquin, as part of a team); 5 Dhaba, tea stall or food stall WORKER (prepared food or tea); 6 Dhaba, tea stall or food stall OWNER (prepared food or tea); 7 Shop WORKER (goods: prasad, puja items, clothes, general store); 8 Shop OWNER (goods, not prepared food); 9 Hotel / lodge WORKER; 10 Hotel / lodge OWNER; 11 Driver (any motor vehicle); 12 Guide; 13 Wage labourer (construction, loading, odd jobs); 14 Other & \texttt{occupation}\newline Source: Project quota design, mapped to NCO-2015 families \\
B2 & In your own words, what exactly is your main work here? Write down what they say — what they actually do, and who for. Do not tick a box for this one. &  & \texttt{occupation\_detail}\newline Source: Project design (coarse quota group + verbatim detail, coded to NCO-2015 in the office) \\
B3 & In this main work, are you: working on your own account without hired workers; running a business with hired workers; a regular monthly wage-earner; or a daily/casual wage-earner? & 1 Own-account (self-employed, no hired workers); 2 Employer/business owner (self-employed); 3 Regular wage or salary; 4 Casual or daily wage labour; 5 Unpaid family worker & \texttt{employment\_type}\newline Source: PLFS / NSS employment status classification [source not held by this project; citation unverified as of 2026-09-28] \\
B4 & Who employs you? A private company; a government or public body; a household (domestic work); a contractor or agent; or an individual employer? & 1 Private company; 2 Government or public body; 3 A household (domestic work); 4 A contractor or agent; 5 An individual employer; 97 Don't know & \texttt{wage\_employer}\newline \textit{Skip: Ask only if employment\_type is 3 or 4 (regular or casual wage worker)}\newline Source: Apablaza et al. 2026, Appendix 2 Q8 (shortened to wage workers); categories ours \\
B5 & Besides your main work, what other paid work do you do during the Yatra season? Check every one that applies; leave blank if there is none. Ask and code the closest group for each; do not read the list aloud unless needed. & 1 Pony/mule worker (works the animals, does not own them); 2 Pony/mule owner (owns the animals); 3 Porter (carries on own back/shoulders - goods, luggage, or a person in a kandi); 4 Palki / dandi bearer (carries a person on a palanquin, as part of a team); 5 Dhaba, tea stall or food stall WORKER (prepared food or tea); 6 Dhaba, tea stall or food stall OWNER (prepared food or tea); 7 Shop WORKER (goods: prasad, puja items, clothes, general store); 8 Shop OWNER (goods, not prepared food); 9 Hotel / lodge WORKER; 10 Hotel / lodge OWNER; 11 Driver (any motor vehicle); 12 Guide; 13 Wage labourer (construction, loading, odd jobs); 14 Other & \texttt{other\_activity\_types}\newline \textit{Skip: May be left blank}\newline Source: Project design (multiple job-holding, common in the Yatra economy) \\
B6 & Taking all of that other work together, about how much do you usually earn from it in a month? (Rs; leave blank if there is no other work) & Rupees (whole number) & \texttt{other\_activity\_income\_pm}\newline \textit{Skip: May be left blank}\newline Source: Project design (multiple job-holding) \\
B7 & Does your household own a shop or stall that you use to earn? & 1 Yes, 0 No & \texttt{owns\_shop\_stall}\newline Source: Chaudhuri et al. 2002 asset covariate \\
B8 & During one Yatra season, how many different employers or businesses do you work in? (1 if you stay with the same one, or run the same one, all season) & Whole number & \texttt{employers\_in\_season}\newline Source: Apablaza et al. 2026 Appendix 2 Q7 (concept: job stability within the engagement); asked as a count per this project's no-scale rule \\
B9 & How many years have you been in this work? Write it as a decimal if needed, for example 1.5 for one year and six months, or 0.5 for six months. & Number & \texttt{years\_current\_job}\newline Source: Apablaza et al. 2026, Appendix 2 Q10; Table 3 (stability) \\
B10 & In the Yatra season, on a normal working day, how many hours do you work in total, across all your work? & Whole number & \texttt{hours\_day\_yatra}\newline Source: Apablaza et al. 2026 Q13 (decomposed into hours/day x days/week rather than asked as one weekly figure, for ease of recall) \\
B11 & How many days a week do you usually work in the season? & Whole number & \texttt{days\_week\_yatra}\newline Source: Apablaza et al. 2026 Q13 (decomposed into hours/day x days/week rather than asked as one weekly figure, for ease of recall) \\
B12 & In the last ten years, did you change your main kind of work? & 1 Yes, 0 No & \texttt{prev\_occ\_change}\newline Source: Occupational mobility check; job-history item \\
B13 & What was your previous main work? Write down exactly what they say, in their words — do not pick from a list. &  & \texttt{prev\_occ}\newline \textit{Skip: Ask only if prev\_occ\_change = 1}\newline Source: Gathmann and Schonberg 2010 (observed moves are systematically shorter than random moves — the validation test for a task-distance measure); coded to NCO-2015 in the office \\
B14 & Looking ahead a few years, what work would you like to be doing? Write down exactly what they say, in their words — do not pick from a list, and do not prompt with examples. If they want to carry on exactly as they are, write that. &  & \texttt{target\_occ}\newline Source: Project design (stated destination, for the structural-displacement analysis); coded to NCO-2015 in the office \\
B15 & What was the main reason you changed? & 1 Better income; 2 Lost the previous work; 3 Work ended with the season; 4 Family reasons; 5 Health or injury; 6 Other & \texttt{prev\_occ\_reason}\newline \textit{Skip: Ask only if prev\_occ\_change = 1}\newline Source: Job-history item \\
\bottomrule
\end{longtable}
}
\subsection*{Module C: Monthly calendar of work and income}
\begin{quote}\small\textit{Read aloud: Now I want to go through the last twelve months, one month at a time. For each month I will ask what work you were mainly doing, and then what you usually earned. It is easiest to start with the Yatra months and fill the rest afterwards.}\end{quote}
{\scriptsize
\begin{longtable}{p{0.8cm}p{6.4cm}p{3.6cm}p{5.8cm}}
\toprule
\textbf{No.} & \textbf{Question or instruction} & \textbf{Answers} & \textbf{Variable, skip rule, source} \\
\midrule
\endfirsthead
\toprule
\textbf{No.} & \textbf{Question or instruction} & \textbf{Answers} & \textbf{Variable, skip rule, source} \\
\midrule
\endhead
C1 & In the past year, how many months did you work here in the Yatra season? Count the months you spent getting ready for it as well. & Whole number & \texttt{months\_worked\_yatra}\newline Source: Project design (replaces the twelve-cell status calendar, 2026-10-05) \\
C2 & And in the past year, how many months did you have no paid work at all? & Whole number & \texttt{months\_no\_paid\_work}\newline Source: Project design (replaces the twelve-cell status calendar, 2026-10-05) \\
C3 & In the remaining months — not the Yatra season, and not the months with no work — what was your main work? Code the closest one. & 2 Farming (crops); 3 Livestock (animals); 4 Casual or daily wage labour; 5 Construction work; 6 Own small shop, stall or trade; 7 Salaried job; 96 Some other work; 0 There were no such months & \texttt{offseason\_activity}\newline \textit{Skip: May be left blank if the Yatra months and the idle months already account for all twelve}\newline Source: Project design (replaces the modal non-Yatra calendar cell, 2026-10-05) \\
C4 & Thinking of the whole past year, about how much did you earn in total from all your work, after costs? (Rs) & Rupees (whole number) & \texttt{income\_annual\_total}\newline \textit{Skip: May be left blank if the respondent declines to discuss earnings}\newline Source: Apablaza et al. 2026 Q15 (annual-total fallback), adapted to a seasonal workforce and now the only earnings route \\
C5 & Out of every 100 rupees of that, how many came from your Yatra work? (The rest is counted as coming from your other work.) & Whole number & \texttt{pct\_income\_yatra}\newline \textit{Skip: Ask only if income\_annual\_total was given; may be left blank}\newline Source: Project design (the Yatra / non-Yatra split needed to reweight an annual total onto the season) \\
C6 & Which range is closest to what you usually earn in a month from work? Under Rs 5,000; Rs 5,000 to 10,000; Rs 10,000 to 20,000; Rs 20,000 or more; or don't know? & 1 Under Rs 5,000; 2 Rs 5,000 to 10,000; 3 Rs 10,000 to 20,000; 4 Rs 20,000 or more; 97 Don't know & \texttt{earnings\_range}\newline \textit{Skip: Ask only if income\_annual\_total was left blank; may be left blank}\newline Source: Apablaza et al. 2026, Appendix 2 Q15; bands to check \\
C7 & In the months when the Yatra is closed and you are doing other work, on a normal working day, how many hours do you work in total? & Whole number & \texttt{hours\_day\_offseason}\newline \textit{Skip: Ask only if there are months that are neither Yatra nor idle (12 - months\_worked\_yatra - months\_no\_paid\_work > 0)}\newline Source: Apablaza et al. 2026 Q13, asked a second time for the off-season (decomposed as hours/day x days/week, as in the Yatra-season pair) \\
C8 & In those months, how many days a week do you usually work? & Whole number & \texttt{days\_week\_offseason}\newline \textit{Skip: Ask only if there are months that are neither Yatra nor idle (12 - months\_worked\_yatra - months\_no\_paid\_work > 0)}\newline Source: Apablaza et al. 2026 Q13, asked a second time for the off-season \\
\bottomrule
\end{longtable}
}
\subsection*{Module D: Migration, home place and remittances}
\begin{quote}\small\textit{Read aloud: Now a few questions about where your home is, and what you do when the Yatra closes for the season.}\end{quote}
{\scriptsize
\begin{longtable}{p{0.8cm}p{6.4cm}p{3.6cm}p{5.8cm}}
\toprule
\textbf{No.} & \textbf{Question or instruction} & \textbf{Answers} & \textbf{Variable, skip rule, source} \\
\midrule
\endfirsthead
\toprule
\textbf{No.} & \textbf{Question or instruction} & \textbf{Answers} & \textbf{Variable, skip rule, source} \\
\midrule
\endhead
D1 & When the Yatra closes for the season, do YOU go to your home place? & 1 Yes, 0 No & \texttt{resp\_returns\_at\_closure}\newline Source: Project design; categories from the pilot's own residency\_pattern distribution (n=46) \\
D2 & For most of the year, do the rest of your household live in your home village or town? & 1 Yes, 0 No & \texttt{hh\_at\_home\_place}\newline Source: Project design (split-household measurement, underpinning Module E) \\
D3 & Is the money you send home different in the Yatra season from when the Yatra is closed? & 1 Yes, 0 No & \texttt{remit\_differs\_by\_season}\newline Source: Project design (mirrors spend\_differs\_by\_season, Module E) \\
D4 & In a normal month during the Yatra season, how much money do you usually send to family or others living elsewhere? (Rs, 0 if none) & Rupees (whole number) & \texttt{remit\_out\_yatra\_pm}\newline Source: Project design (amount only; frequency dropped for length). A previous citation to NSS 64th Round practice was withdrawn 2026-09-28: no NSS schedule is held by this project and it could not be checked \\
D5 & In a normal month when the Yatra is closed, how much do you usually send? (Rs, 0 if none) & Rupees (whole number) & \texttt{remit\_out\_offseason\_pm}\newline \textit{Skip: Ask only if remit\_differs\_by\_season = 1}\newline Source: Project design (amount only; frequency dropped for length). A previous citation to NSS 64th Round practice was withdrawn 2026-09-28: no NSS schedule is held by this project and it could not be checked \\
D6 & How do you usually send it? & 1 From my own phone (UPI or net banking); 2 In person at a bank branch or bank mitra / CSP; 3 Money order or post office; 4 Sent with someone going home; 5 Carried it myself; 6 Through an agent or middleman; 7 Other & \texttt{remit\_mode}\newline \textit{Skip: Ask only if remit\_out\_yatra\_pm > 0 or remit\_out\_offseason\_pm > 0}\newline Source: Project design (remittance channel as a financial-inclusion and cost measure) \\
D7 & In a normal month during the Yatra season, how much money do you usually receive from family or others? (Rs, 0 if none) & Rupees (whole number) & \texttt{remit\_in\_yatra\_pm}\newline Source: Project design (amount only; frequency dropped for length). A previous citation to NSS 64th Round practice was withdrawn 2026-09-28: no NSS schedule is held by this project and it could not be checked \\
D8 & In a normal month when the Yatra is closed, how much do you usually receive? (Rs, 0 if none) & Rupees (whole number) & \texttt{remit\_in\_offseason\_pm}\newline \textit{Skip: Ask only if remit\_differs\_by\_season = 1}\newline Source: Project design (amount only; frequency dropped for length). A previous citation to NSS 64th Round practice was withdrawn 2026-09-28: no NSS schedule is held by this project and it could not be checked \\
D9 & How many people from your household work in the same business or establishment as you, counting yourself? & Whole number & \texttt{n\_hh\_same\_business}\newline Source: Project design (enterprise concentration of household labour, 2026-10-05) \\
D10 & How many years have you been coming here for the Yatra season? & Whole number & \texttt{years\_coming\_here}\newline \textit{Skip: Ask only if resp\_returns\_at\_closure = 1}\newline Source: Project design (duration of the seasonal relationship) \\
D11 & What was the main reason you first came here to work? & 1 No work at home; 2 Pay is better here; 3 Family or people from my village were already here; 4 Land at home is too little to live on; 5 Debt to repay; 6 A contractor or agent brought me; 7 Married into / moved with family; 8 Some other reason & \texttt{came\_here\_reason}\newline \textit{Skip: Ask only if origin is not Local (same district)}\newline Source: Project design (push/pull as a vulnerability covariate) \\
D12 & What was that reason? Write down what they say. &  & \texttt{came\_here\_reason\_other}\newline \textit{Skip: Ask only if came\_here\_reason = 8 (some other reason). May be left blank.}\newline Source: Project design (catch-all codes lose the case they catch) \\
D13 & Last year, in the months when the Yatra was closed, did you work away from your home place? & 1 Yes, 0 No & \texttt{worked\_away\_in\_closure}\newline Source: Project design (off-season labour migration) \\
D14 & Where did you go, and what work did you do there? Write what they say. &  & \texttt{closure\_work\_detail}\newline \textit{Skip: Ask only if worked\_away\_in\_closure = 1; may be left blank}\newline Source: Project design; coded to NCO-2015 in the office \\
D15 & If this work here ended, would you go away from your home place to look for work in the coming year? & 0 No; 1 Yes; 97 Don't know & \texttt{would\_move\_for\_work}\newline Source: Project design; horizoned intention item after VASyR 2025 (move\_accom\_yesno, asked over a stated 6-month horizon) \\
D16 & Who mainly helped you get this work, or arranged it for you? & 1 A family member already working here; 2 A friend or someone from the village; 3 A thekedar or contractor I now work for; 4 An agent or middleman who placed me (usually for a fee); 5 No one; I found it myself; 6 The employer called me directly; 7 Someone else & \texttt{migration\_referral}\newline Source: Project design (referral channel as a proxy for social capital and for placement debt) \\
D17 & Who was it? Write what they say. &  & \texttt{migration\_referral\_other}\newline \textit{Skip: Ask only if migration\_referral = 7 (Other); may be left blank}\newline Source: Standard household-survey item \\
\bottomrule
\end{longtable}
}
\subsection*{Module E: Household consumption}
\begin{quote}\small\textit{Read aloud: Now about what your household usually spends in a month. For each thing I will ask twice — once for a normal month during the Yatra season, and once for a normal month when the Yatra is closed — because spending is often quite different in the two. Rough figures are fine.}\end{quote}
{\scriptsize
\begin{longtable}{p{0.8cm}p{6.4cm}p{3.6cm}p{5.8cm}}
\toprule
\textbf{No.} & \textbf{Question or instruction} & \textbf{Answers} & \textbf{Variable, skip rule, source} \\
\midrule
\endfirsthead
\toprule
\textbf{No.} & \textbf{Question or instruction} & \textbf{Answers} & \textbf{Variable, skip rule, source} \\
\midrule
\endhead
E1 & Leaving aside money you send home — is what your household spends in a normal month during the Yatra season different from what it spends when the Yatra is closed? & 1 Yes, 0 No & \texttt{spend\_differs\_by\_season}\newline Source: Project design (respondent-gated seasonal recall); the underlying two-season design follows Beegle et al. 2012 and Deaton and Grosh 2000 on usual-period recall \\
E2 & In a normal month during the Yatra season, how much does your household spend on cereals (rice, wheat, other grains), pulses, sugar and salt, bought or from your own stock? (Rs) & Rupees (whole number) & \texttt{cons\_staples\_yatra\_pm}\newline Source: HCES 2022-23, MoSPI Appendix A, Sections 5.1-5.3 (30-day); IHDS-II Income and Social Capital Questionnaire, Q14 14.1-14.6 (30-day, same items) \\
E3 & In a normal month when the Yatra is closed, how much does your household spend on cereals (rice, wheat, other grains), pulses, sugar and salt, bought or from your own stock? (Rs) & Rupees (whole number) & \texttt{cons\_staples\_offseason\_pm}\newline \textit{Skip: Ask only if spend\_differs\_by\_season = 1}\newline Source: HCES 2022-23, MoSPI Appendix A, Sections 5.1-5.3 (30-day); IHDS-II Income and Social Capital Questionnaire, Q14 14.1-14.6 (30-day, same items) \\
E4 & In a normal month during the Yatra season, how much does your household spend on milk and milk products, vegetables, fruit, egg/fish/meat, cooking oil, spices, and tea or coffee, bought or from your own stock? (Rs) & Rupees (whole number) & \texttt{cons\_perishables\_yatra\_pm}\newline Source: HCES 2022-23, MoSPI Appendix A, Sections 6.1-6.8 (7-day actual recall in HCES; asked here as a usual monthly figure instead, for the reason above) \\
E5 & In a normal month when the Yatra is closed, how much does your household spend on milk and milk products, vegetables, fruit, egg/fish/meat, cooking oil, spices, and tea or coffee, bought or from your own stock? (Rs) & Rupees (whole number) & \texttt{cons\_perishables\_offseason\_pm}\newline \textit{Skip: Ask only if spend\_differs\_by\_season = 1}\newline Source: HCES 2022-23, MoSPI Appendix A, Sections 6.1-6.8 (7-day actual recall in HCES; asked here as a usual monthly figure instead, for the reason above) \\
E6 & If your household had not grown, raised or been given any of your food, about how much would it have cost to buy in a normal month during the Yatra season? (Rs, 0 if you buy everything) & Rupees (whole number) & \texttt{cons\_food\_own\_yatra\_pm}\newline Source: Calvo and Dercon 2007; Eze and Iheonu 2025 (non-purchased food, aggregate approach) \\
E7 & If your household had not grown, raised or been given any of your food, about how much would it have cost to buy in a normal month when the Yatra is closed? (Rs, 0 if you buy everything) & Rupees (whole number) & \texttt{cons\_food\_own\_offseason\_pm}\newline \textit{Skip: Ask only if spend\_differs\_by\_season = 1}\newline Source: Calvo and Dercon 2007; Eze and Iheonu 2025 (non-purchased food, aggregate approach) \\
E8 & In a normal month during the Yatra season, how much does your household spend on meals, tea and snacks eaten outside the home, including any an employer gave free, at their market value? (Rs) & Rupees (whole number) & \texttt{cons\_food\_out\_yatra\_pm}\newline Source: HCES 2022-23, MoSPI Appendix A, Section 7.1 'served processed food' (7-day in HCES; usual monthly here) \\
E9 & In a normal month when the Yatra is closed, how much does your household spend on meals, tea and snacks eaten outside the home, including any an employer gave free, at their market value? (Rs) & Rupees (whole number) & \texttt{cons\_food\_out\_offseason\_pm}\newline \textit{Skip: Ask only if spend\_differs\_by\_season = 1}\newline Source: HCES 2022-23, MoSPI Appendix A, Section 7.1 'served processed food' (7-day in HCES; usual monthly here) \\
E10 & In a normal month during the Yatra season, how much does your household spend on fuel and light: cooking fuel, firewood, electricity, kerosene, candles? (Rs) & Rupees (whole number) & \texttt{cons\_fuel\_yatra\_pm}\newline Source: HCES 2022-23, MoSPI Appendix A, Section 8.1 (30-day). IHDS-II Income and Social Capital Questionnaire, Q14 14.21-14.22. Nigeria GHS-Panel Wave 3, Household Questionnaire, Sections 10B/11 (30-day tier) \\
E11 & In a normal month when the Yatra is closed, how much does your household spend on fuel and light: cooking fuel, firewood, electricity, kerosene, candles? (Rs) & Rupees (whole number) & \texttt{cons\_fuel\_offseason\_pm}\newline \textit{Skip: Ask only if spend\_differs\_by\_season = 1}\newline Source: HCES 2022-23, MoSPI Appendix A, Section 8.1 (30-day). IHDS-II Income and Social Capital Questionnaire, Q14 14.21-14.22. Nigeria GHS-Panel Wave 3, Household Questionnaire, Sections 10B/11 (30-day tier) \\
E12 & In a normal month during the Yatra season, how much does your household spend on soap, toiletries, cleaning goods and other small household items? (Rs) & Rupees (whole number) & \texttt{cons\_routine\_misc\_yatra\_pm}\newline Source: HCES 2022-23, MoSPI Appendix A, Sections 9.1-9.2 (30-day). IHDS-II Income and Social Capital Questionnaire, Q14 14.25-14.27. Nigeria GHS-Panel Wave 3, Household Questionnaire, Sections 10B/11 (30-day tier) \\
E13 & In a normal month when the Yatra is closed, how much does your household spend on soap, toiletries, cleaning goods and other small household items? (Rs) & Rupees (whole number) & \texttt{cons\_routine\_misc\_offseason\_pm}\newline \textit{Skip: Ask only if spend\_differs\_by\_season = 1}\newline Source: HCES 2022-23, MoSPI Appendix A, Sections 9.1-9.2 (30-day). IHDS-II Income and Social Capital Questionnaire, Q14 14.25-14.27. Nigeria GHS-Panel Wave 3, Household Questionnaire, Sections 10B/11 (30-day tier) \\
E14 & In a normal month during the Yatra season, how much does your household spend on local transport (bus, shared jeep, auto) and phone or internet charges? (Rs) & Rupees (whole number) & \texttt{cons\_transport\_comm\_yatra\_pm}\newline Source: HCES 2022-23, MoSPI Appendix A, Sections 11.1-11.2 (30-day). IHDS-II Income and Social Capital Questionnaire, Q14 14.24/14.28. Nigeria GHS-Panel Wave 3, Household Questionnaire, Sections 10B/11 (30-day tier) \\
E15 & In a normal month when the Yatra is closed, how much does your household spend on local transport (bus, shared jeep, auto) and phone or internet charges? (Rs) & Rupees (whole number) & \texttt{cons\_transport\_comm\_offseason\_pm}\newline \textit{Skip: Ask only if spend\_differs\_by\_season = 1}\newline Source: HCES 2022-23, MoSPI Appendix A, Sections 11.1-11.2 (30-day). IHDS-II Income and Social Capital Questionnaire, Q14 14.24/14.28. Nigeria GHS-Panel Wave 3, Household Questionnaire, Sections 10B/11 (30-day tier) \\
E16 & In a normal month during the Yatra season, how much does your household spend on rent for the home (0 if owned, and 0 for any period spent living for free)? (Rs) & Rupees (whole number) & \texttt{cons\_rent\_yatra\_pm}\newline Source: HCES 2022-23, MoSPI Appendix A, Section 11.4 (30-day). IHDS-II Income and Social Capital Questionnaire, Q14 14.30. Nigeria GHS-Panel Wave 3, Household Questionnaire, Sections 10B/11 (30-day tier) \\
E17 & In a normal month when the Yatra is closed, how much does your household spend on rent for the home (0 if owned, and 0 for any period spent living for free)? (Rs) & Rupees (whole number) & \texttt{cons\_rent\_offseason\_pm}\newline \textit{Skip: Ask only if spend\_differs\_by\_season = 1}\newline Source: HCES 2022-23, MoSPI Appendix A, Section 11.4 (30-day). IHDS-II Income and Social Capital Questionnaire, Q14 14.30. Nigeria GHS-Panel Wave 3, Household Questionnaire, Sections 10B/11 (30-day tier) \\
E18 & In a normal month during the Yatra season, how much does your household spend on medicine, doctor's or clinic fees and tests, not counting a hospital stay? (Rs) & Rupees (whole number) & \texttt{cons\_med\_nonhosp\_yatra\_pm}\newline Source: HCES 2022-23, MoSPI Appendix A, Section 10.3 (30-day). IHDS-II Income and Social Capital Questionnaire, Q14 14.33 (30-day, out-patient) \\
E19 & In a normal month when the Yatra is closed, how much does your household spend on medicine, doctor's or clinic fees and tests, not counting a hospital stay? (Rs) & Rupees (whole number) & \texttt{cons\_med\_nonhosp\_offseason\_pm}\newline \textit{Skip: Ask only if spend\_differs\_by\_season = 1}\newline Source: HCES 2022-23, MoSPI Appendix A, Section 10.3 (30-day). IHDS-II Income and Social Capital Questionnaire, Q14 14.33 (30-day, out-patient) \\
E20 & In a normal month during the Yatra season, how much does your household spend on biscuits, namkeen, chips, packaged snacks, cold drinks or bottled water? (Rs) & Rupees (whole number) & \texttt{cons\_packaged\_food\_yatra\_pm}\newline Source: HCES 2022-23, MoSPI Appendix A, Section 7.2 (7-day in HCES; usual monthly here, as for every other seasonal item) \\
E21 & In a normal month when the Yatra is closed, how much does your household spend on biscuits, namkeen, chips, packaged snacks, cold drinks or bottled water? (Rs) & Rupees (whole number) & \texttt{cons\_packaged\_food\_offseason\_pm}\newline \textit{Skip: Ask only if spend\_differs\_by\_season = 1}\newline Source: HCES 2022-23, MoSPI Appendix A, Section 7.2 (7-day in HCES; usual monthly here, as for every other seasonal item) \\
E22 & In a normal month during the Yatra season, how much does your household spend on pan, gutka, bidi, cigarettes, tobacco or alcohol? (Rs) & Rupees (whole number) & \texttt{cons\_pan\_tobacco\_yatra\_pm}\newline Source: HCES 2022-23, MoSPI Appendix A, Sections 12.1-12.3 (7-day in HCES; usual monthly here, and combined into one figure) \\
E23 & In a normal month when the Yatra is closed, how much does your household spend on pan, gutka, bidi, cigarettes, tobacco or alcohol? (Rs) & Rupees (whole number) & \texttt{cons\_pan\_tobacco\_offseason\_pm}\newline \textit{Skip: Ask only if spend\_differs\_by\_season = 1}\newline Source: HCES 2022-23, MoSPI Appendix A, Sections 12.1-12.3 (7-day in HCES; usual monthly here, and combined into one figure) \\
E24 & In the last 12 months, how much did your household spend on clothes and footwear? (Rs) & Rupees (whole number) & \texttt{cons\_clothing\_12m}\newline Source: HCES 2022-23, MoSPI Appendix A, Sections 13.1-13.2 (365-day). IHDS-II Income and Social Capital Questionnaire, Q14 14.38-14.39 (365-day). Nigeria GHS-Panel Wave 3, Household Questionnaire, Sections 10B/11 instead uses a 6-month recall for clothing (items 401-441); we follow the two India sources, since one of them (HCES) underlies our poverty line \\
E25 & In the last 12 months, how much on education: fees, books, uniforms, tuition? (Rs) & Rupees (whole number) & \texttt{cons\_education\_12m}\newline Source: HCES 2022-23, MoSPI Appendix A, Section 10.1 (365-day). IHDS-II Income and Social Capital Questionnaire, Q14 14.35-14.37 (365-day, same items). Nigeria GHS-Panel Wave 3, Household Questionnaire, Sections 10B/11 asks school fees inside its education/roster module, per child, not as a household total, so it is not directly comparable \\
E26 & In the last 12 months, how much did your household pay for any hospital admission or stay? (Rs, 0 if none) & Rupees (whole number) & \texttt{cons\_medical\_hosp\_12m}\newline Source: HCES 2022-23, MoSPI Appendix A, Section 10.2 (365-day, hospitalisation). IHDS-II Income and Social Capital Questionnaire, Q14 14.34 (365-day, in-patient). Nigeria GHS-Panel Wave 3, Household Questionnaire, Sections 10B/11 instead uses one combined 6-month medical item (less detail); we follow the matching HCES/IHDS split \\
E27 & In the last 12 months, how much on durable goods: furniture, utensils, cooking appliances, phone, jewellery or ornaments, bicycle or vehicle parts? (Rs) & Rupees (whole number) & \texttt{cons\_durables\_12m}\newline Source: HCES 2022-23, MoSPI Appendix A, Section 14 (365-day, 10 sub-groups). IHDS-II Income and Social Capital Questionnaire, Q14 14.40-14.48 (365-day). Nigeria GHS-Panel Wave 3, Household Questionnaire, Sections 10B/11 (12-month tier: durables, appliances, building materials). All three agree on the 12-month recall; insurance premiums, which IHDS (14.50) and the Nigeria survey (items 510-513) count here, are excluded, following HCES and standard national-accounts practice (insurance is a financial item, not consumption) — premiums are covered instead by the insurance items in Module G \\
\bottomrule
\end{longtable}
}
\subsection*{Module F: Housing, amenities and assets}
\begin{quote}\small\textit{Read aloud: Now some questions about your house and the things your household has.}\end{quote}
{\scriptsize
\begin{longtable}{p{0.8cm}p{6.4cm}p{3.6cm}p{5.8cm}}
\toprule
\textbf{No.} & \textbf{Question or instruction} & \textbf{Answers} & \textbf{Variable, skip rule, source} \\
\midrule
\endfirsthead
\toprule
\textbf{No.} & \textbf{Question or instruction} & \textbf{Answers} & \textbf{Variable, skip rule, source} \\
\midrule
\endhead
F1 & What is the main material of the floor at your usual home? & 1 Mud/kaccha; 2 Cement/mud-cement; 3 Tile/mosaic/marble & \texttt{floor\_material}\newline Source: NITI Aayog National MPI housing indicator \\
F2 & What is the main material of the roof at your usual home? & 1 Thatch/wood/mud; 2 Tin/GI sheet; 3 Concrete/RCC & \texttt{roof\_material}\newline Source: NITI Aayog National MPI housing indicator \\
F3 & What are the walls of your usual home mainly made of? & 1 Mud, thatch, bamboo or other natural material; 2 Unburnt brick, wood or tin; 3 Burnt brick, cement, concrete or stone & \texttt{wall\_material}\newline Source: NITI Aayog National MPI housing indicator (wall limb) \\
F4 & While you are here for the season, where do you sleep? & 1 At my own usual home (I live here); 2 Rented room or house; 3 Room or dormitory the employer provides; 4 Sleeps at the shop, dhaba or workplace; 5 Tent or temporary shelter; 6 Lodge, dharamshala or ashram; 7 In the open, or a verandah; 8 Somewhere else & \texttt{accom\_type\_here}\newline Source: Lyons et al. 2023, Table 2 (area/settlement conditions), adapted to a labour-migrant setting \\
F5 & Does your usual home have an electricity connection? & 1 Yes, 0 No & \texttt{electricity}\newline Source: NITI Aayog National MPI \\
F6 & What kind of toilet does your household use at your usual home? & 1 No toilet / open defecation; 2 Pit latrine without slab or open pit; 3 Improved toilet, but shared with other households; 4 Improved toilet, used only by this household & \texttt{toilet\_type}\newline Source: NITI Aayog National MPI sanitation indicator (improved/unimproved and shared/not) \\
F7 & What is your main source of drinking water at your usual home? & 1 Piped into the house or yard; 2 Public tap or standpipe; 3 Handpump, tubewell or borewell; 4 Protected well or protected spring; 5 Rainwater collection; 6 Bottled, packaged or community RO; 7 Unprotected well or unprotected spring; 8 River, stream, pond or canal; 9 Tanker truck or cart with drum & \texttt{drinking\_water}\newline Source: NITI Aayog National MPI drinking-water indicator (JMP improved/unimproved taxonomy) \\
F8 & Is the drinking water available at the house itself? & 1 Yes, 0 No & \texttt{water\_on\_premises}\newline Source: NITI Aayog National MPI drinking-water indicator (30-minute round-trip rule) \\
F9 & How long does it take to go there, get the water and come back? (minutes, round trip) & Whole number & \texttt{water\_fetch\_minutes}\newline \textit{Skip: Ask only if water\_on\_premises = 0}\newline Source: NITI Aayog National MPI drinking-water indicator (30-minute round-trip rule) \\
F10 & What does your household mainly use to cook at your usual home? & 1 LPG or cylinder gas; 2 Piped natural gas; 3 Electricity; 4 Biogas (gobar gas plant); 5 Kerosene; 6 Firewood; 7 Dung cakes (gobar); 8 Crop residue, straw or shrubs; 9 Charcoal; 10 Coal or lignite & \texttt{cooking\_fuel}\newline Source: NITI Aayog National MPI cooking-fuel indicator (its own list of dirty fuels) \\
F11 & In what unit do you count your land — nali, bigha, acres or hectares? & 1 Nali; 2 Bigha; 3 Acres; 4 Hectares; 5 No land & \texttt{land\_unit}\newline Source: Project design (local land units, Uttarakhand) \\
F12 & And how much CULTIVABLE land is that? Do not count the land the house stands on. & Number & \texttt{land\_cultivable\_acres}\newline Source: Chaudhuri et al. 2002 asset covariate \\
F13 & Which of these does your household own? Read the list out and check every one they say: television, radio, bicycle, motorcycle or scooter, car/jeep/truck, telephone of any kind, computer or laptop, animal cart, refrigerator. & 1 Television; 2 Radio; 3 Bicycle; 4 Motorcycle or scooter; 5 Car, jeep or truck; 6 Telephone of any kind (mobile or landline); 7 Computer or laptop; 8 Cart pulled by an animal; 9 Refrigerator & \texttt{assets\_owned}\newline \textit{Skip: May be left blank if the household owns none of them}\newline Source: NITI Aayog National MPI, Assets (weight 1/21) \\
F14 & And which of these animals does your household own? Cows or buffaloes, goats or sheep, ponies or mules. & 1 Cows or buffaloes; 2 Goats or sheep; 3 Ponies or mules & \texttt{livestock\_owned}\newline \textit{Skip: May be left blank if the household owns no animals}\newline Source: Chaudhuri et al. 2002 asset covariate; pony/mule class is this project's exposure variable \\
F15 & Do you use a vehicle to earn money — for example hauling goods, transporting people, or running a taxi? It does not have to be yours. & 1 Yes, 0 No & \texttt{owns\_work\_vehicle}\newline Source: Chaudhuri et al. 2002 asset covariate \\
F16 & Does your household own equipment or tools that you use to earn? & 1 Yes, 0 No & \texttt{owns\_work\_equipment}\newline Source: Chaudhuri et al. 2002 asset covariate \\
F17 & What are they? Write what they say. Leave blank if they cannot say. &  & \texttt{work\_equipment\_detail}\newline \textit{Skip: Ask only if owns\_work\_equipment = 1; may be left blank}\newline Source: Standard household-survey item \\
\bottomrule
\end{longtable}
}
\subsection*{Module G: Finance, insurance and schemes}
\begin{quote}\small\textit{Read aloud: Now about savings, loans, insurance and government schemes.}\end{quote}
{\scriptsize
\begin{longtable}{p{0.8cm}p{6.4cm}p{3.6cm}p{5.8cm}}
\toprule
\textbf{No.} & \textbf{Question or instruction} & \textbf{Answers} & \textbf{Variable, skip rule, source} \\
\midrule
\endfirsthead
\toprule
\textbf{No.} & \textbf{Question or instruction} & \textbf{Answers} & \textbf{Variable, skip rule, source} \\
\midrule
\endhead
G1 & Does anyone in your household have a bank account or a post office account? & 1 Yes, 0 No & \texttt{has\_bank\_account}\newline Source: NITI Aayog National MPI bank-account indicator; VEP adaptive-capacity covariate \\
G2 & In the last 12 months, did you or anyone in your household borrow money? & 1 Yes, 0 No & \texttt{took\_loan\_12m}\newline Source: VEP adaptive-capacity covariate \\
G3 & Who was the main lender? & 1 Nationalised or public-sector bank; 2 Private bank; 3 Cooperative bank or RRB; 4 Microfinance institution or SHG; 5 Moneylender; 6 Relative or friend; 7 Employer or contractor (advance); 8 Other & \texttt{credit\_source}\newline \textit{Skip: Ask only if took\_loan\_12m = 1}\newline Source: VEP adaptive-capacity covariate (lender type) \\
G4 & What did you take it for? & 1 Medical or hospital costs; 2 Buying an animal, vehicle or equipment for work; 3 Shop stock or business; 4 Food or daily household needs; 5 A wedding, funeral or ceremony; 6 House building or repair; 7 Education; 8 Repaying another loan; 9 Something else & \texttt{loan\_purpose}\newline \textit{Skip: Ask only if took\_loan\_12m = 1}\newline Source: Project design (investment against distress borrowing) \\
G5 & How much did you borrow in total? (Rs) & Rupees (whole number) & \texttt{loan\_amount\_borrowed}\newline \textit{Skip: Ask only if took\_loan\_12m = 1}\newline Source: Project design (loan size, for the debt-burden measures) \\
G6 & How much of that loan is still left to repay? (Rs) & Rupees (whole number) & \texttt{loan\_amount}\newline \textit{Skip: Ask only if took\_loan\_12m = 1}\newline Source: Islam and Chowdhury 2025 (financial distress and household vulnerability to poverty) \\
G7 & Do you pay interest on it? & 1 Yes, 0 No & \texttt{pays\_interest}\newline \textit{Skip: Ask only if took\_loan\_12m = 1}\newline Source: Project design (informal lending is often interest-free) \\
G8 & On every 100 rupees you borrowed, about how much interest do you pay in a month? (Rs; 0 if none, leave blank if not known) & Number & \texttt{loan\_interest\_per100\_pm}\newline \textit{Skip: Ask only if took\_loan\_12m = 1; may be left blank}\newline Source: Project design (informal-credit pricing as locally quoted) \\
G9 & Did you have to give anything as security for it? If so, what — jewellery, land, animals, a vehicle, shop stock, a house, or something else? & 0 Nothing was given as security; 1 Jewellery or ornaments; 2 Land; 3 Animals; 4 A vehicle; 5 Shop stock or business goods; 6 House or building; 7 Something else & \texttt{loan\_collateral}\newline \textit{Skip: Ask only if took\_loan\_12m = 1}\newline Source: Carter and Zimmerman 2000; Zimmerman and Carter 2003 (asset-based coping) \\
G10 & How many people in your household are covered by any health insurance or health scheme, including government ones? (0 if none) & Whole number & \texttt{n\_health\_insured}\newline Source: NITI Aayog National MPI health-insurance indicator; Lyons et al. 2023 \\
G11 & How many people in your household have life insurance? (0 if none) & Whole number & \texttt{n\_life\_insured}\newline Source: Standard household-survey item \\
G12 & Is your crop insured? & 1 Yes, 0 No & \texttt{has\_crop\_insurance}\newline \textit{Skip: Ask only if land\_cultivable\_acres > 0}\newline Source: Standard household-survey item \\
G13 & In the last 12 months, did your household get anything from any government scheme? Read the list out before you take an answer. & 1 Yes, 0 No & \texttt{govt\_any\_benefit}\newline Source: Project design; scheme names as administered in Uttarakhand \\
G14 & Which ones? Write down every scheme they name, in their own words. &  & \texttt{govt\_schemes\_detail}\newline \textit{Skip: Ask only if govt\_any\_benefit = 1}\newline Source: Project design; coded in the office against the Uttarakhand scheme list \\
G15 & Do you own a smartphone? & 1 Yes, 0 No & \texttt{smartphone\_owned}\newline Source: VEP adaptive-capacity covariate; STEP digital items \\
G16 & How many people in your household can make a payment by phone themselves, without help? (0 if none) & Whole number & \texttt{n\_can\_transact\_online}\newline Source: Project design (intra-household financial capability) \\
\bottomrule
\end{longtable}
}
\subsection*{Module H: Health}
\begin{quote}\small\textit{Read aloud: Now a few questions about health in your household. You may leave out any of these.}\end{quote}
{\scriptsize
\begin{longtable}{p{0.8cm}p{6.4cm}p{3.6cm}p{5.8cm}}
\toprule
\textbf{No.} & \textbf{Question or instruction} & \textbf{Answers} & \textbf{Variable, skip rule, source} \\
\midrule
\endfirsthead
\toprule
\textbf{No.} & \textbf{Question or instruction} & \textbf{Answers} & \textbf{Variable, skip rule, source} \\
\midrule
\endhead
H1 & In the last 15 days, was anyone in your household ill? & 1 Yes, 0 No & \texttt{morbidity\_15d}\newline \textit{Skip: May be left blank}\newline Source: NSS health module (15-day recall) [source not held by this project; citation unverified as of 2026-09-28] \\
H2 & How did the household mainly cope with the cost of this: used savings; borrowed money; sold or pawned assets; cut other consumption; got help from relatives or community; or something else? & 1 Used savings; 2 Borrowed money; 3 Sold or pawned assets; 4 Cut consumption; 5 Help from relatives/community; 6 Paid it out of normal earnings, nothing given up; 7 Something else (say what) & \texttt{morbidity\_coping\_15d}\newline \textit{Skip: Ask only if morbidity\_15d = 1. May be left blank}\newline Source: Project design (a reported illness with no cost/coping follow-up said nothing about its burden) \\
H3 & About how much did the household spend on this in the last 15 days (medicine, doctor or clinic fees, travel for care)? (Rs) & Rupees (whole number) & \texttt{morbidity\_cost\_15d}\newline \textit{Skip: Ask only if morbidity\_15d = 1. May be left blank}\newline Source: Project design \\
H4 & Do you have difficulty walking or climbing steps? & 1 No difficulty; 2 Some difficulty; 3 A lot of difficulty; 4 Cannot do it at all & \texttt{func\_limitation}\newline \textit{Skip: May be left blank}\newline Source: Washington Group on Disability Statistics, Short Set item 3 (mobility), verbatim; Lyons et al. 2023 health dimension \\
H5 & In the last 12 months, was anyone in your household admitted to hospital overnight? & 1 Yes, 0 No & \texttt{hospitalization\_365d}\newline \textit{Skip: May be left blank}\newline Source: NSS health module (365-day recall) [source not held by this project; citation unverified as of 2026-09-28] \\
H6 & In the last 3 months, was there a time when someone in your household needed medical care but could not get it? & 1 Yes, 0 No & \texttt{health\_access\_barrier\_3m}\newline \textit{Skip: May be left blank}\newline Source: VASyR 2025 barriers\_health\_case\_access\_phc\_m (primary-care access barriers); Lyons et al. 2023 Table 2 indicator 2 'healthcare access' \\
H7 & In the last five years, has any child or young person under 18 in your household died? & 1 Yes, 0 No & \texttt{child\_death\_5y}\newline \textit{Skip: May be left blank}\newline Source: NITI Aayog National MPI, Child and Adolescent Mortality (weight 1/12) \\
H8 & In the last five years, did any woman in your household give birth? & 1 Yes, 0 No & \texttt{birth\_last\_5y}\newline \textit{Skip: May be left blank}\newline Source: NITI Aayog National MPI, Maternal Health (weight 1/12) \\
H9 & For the most recent birth, did she have at least four check-ups before the delivery? & 0 No; 1 Yes & \texttt{anc\_4\_visits}\newline \textit{Skip: Ask only if birth\_last\_5y = 1. May be left blank}\newline Source: NITI Aayog National MPI, Maternal Health (antenatal care limb) \\
H10 & Who conducted that delivery — a doctor, a nurse or ANM, an ASHA or Anganwadi worker, a dai, or nobody trained? & 1 Doctor; 2 Nurse, ANM or LHV; 3 ASHA or Anganwadi worker only; 4 Dai (traditional birth attendant); 5 Nobody trained; family only & \texttt{skilled\_birth\_attendant}\newline \textit{Skip: Ask only if birth\_last\_5y = 1. May be left blank}\newline Source: NITI Aayog National MPI, Maternal Health (assisted delivery limb); skilled-provider definition per NFHS \\
\bottomrule
\end{longtable}
}
\subsection*{Module I: Food security, shocks and coping}
\begin{quote}\small\textit{Read aloud: Now about food, and about any difficult times in the last year and how your household managed.}\end{quote}
{\scriptsize
\begin{longtable}{p{0.8cm}p{6.4cm}p{3.6cm}p{5.8cm}}
\toprule
\textbf{No.} & \textbf{Question or instruction} & \textbf{Answers} & \textbf{Variable, skip rule, source} \\
\midrule
\endfirsthead
\toprule
\textbf{No.} & \textbf{Question or instruction} & \textbf{Answers} & \textbf{Variable, skip rule, source} \\
\midrule
\endhead
I1 & In the last 12 months, did any of these happen — to your household, or to the Yatra route you work on? Check every one that applies. & 9 Nothing of this kind happened; 1 Serious illness or death of an earning member; 2 Crop failure or livestock loss; 3 Damage from a natural disaster (flood, landslide, fire); 4 Loss of business, goods or assets; 5 Road or bridge blocked, route closed; 6 Yatra stopped early or badly disrupted; 7 Sharp fall in customers or prices; 8 Lost the work or the work stopped & \texttt{distress\_event\_last365d}\newline Source: VEP exposure covariate (Azeem et al. 2016); shock inventory following Gunther and Harttgen 2009, via Fujii 2016 section 4 \\
I2 & Thinking of whichever of those hit you hardest — about how many weeks of work did you lose because of it? (0 if none) & Whole number & \texttt{shock\_work\_lost\_weeks}\newline \textit{Skip: Ask only if distress\_event\_last365d names a real shock (not code 9, nothing happened)}\newline Source: Ligon and Schechter 2003; Dercon and Krishnan 2000 (VER needs shock magnitude, not only incidence) \\
I3 & And about how much money did your household have to spend because of it — treatment, repairs, replacing what was lost? (Rs, 0 if nothing) & Rupees (whole number) & \texttt{shock\_money\_spent}\newline \textit{Skip: Ask only if distress\_event\_last365d names a real shock (not code 9, nothing happened)}\newline Source: Ligon and Schechter 2003; Dercon and Krishnan 2000 (VER needs shock magnitude, not only incidence) \\
I4 & And in which month did that one happen? & 1 January; 2 February; 3 March; 4 April; 5 May; 6 June; 7 July; 8 August; 9 September; 10 October; 11 November; 12 December & \texttt{shock\_month}\newline \textit{Skip: Ask only if distress\_event\_last365d names a real shock (not code 9, nothing happened)}\newline Source: Dercon and Krishnan 2000 (seasonal timing of shocks); project design \\
I5 & How did your household manage? Check every one they used. & 1 Used savings; 2 Borrowed money; 3 Sold or pawned assets; 4 Cut consumption; 5 Help from relatives/community; 6 Paid it out of normal earnings, nothing given up; 7 Something else (say what) & \texttt{shock\_coping}\newline \textit{Skip: Ask only if any shock was reported}\newline Source: Coping strategies in VEP studies; Carter and Zimmerman 2000 on asset versus consumption smoothing \\
I6 & In the last 12 months, was there a time when, because of a lack of money or other resources, you had to skip a meal? & 1 Yes, 0 No & \texttt{fies\_skipped\_meal}\newline \textit{Skip: May be left blank}\newline Source: FAO Food Insecurity Experience Scale (FIES), SDG indicator 2.1.2; behavioural items only \\
I7 & In the last 12 months, was there a time when, because of a lack of money or other resources, you ate less than you thought you should? & 1 Yes, 0 No & \texttt{fies\_ate\_less}\newline \textit{Skip: May be left blank}\newline Source: FAO Food Insecurity Experience Scale (FIES), SDG indicator 2.1.2; behavioural items only \\
I8 & In the last 12 months, was there a time when, because of a lack of money or other resources, your household ran out of food? & 1 Yes, 0 No & \texttt{fies\_ran\_out}\newline \textit{Skip: May be left blank}\newline Source: FAO Food Insecurity Experience Scale (FIES), SDG indicator 2.1.2; behavioural items only \\
I9 & In the last 12 months, was there a time when, because of a lack of money or other resources, you were hungry but did not eat? & 1 Yes, 0 No & \texttt{fies\_hungry}\newline \textit{Skip: May be left blank}\newline Source: FAO Food Insecurity Experience Scale (FIES), SDG indicator 2.1.2; behavioural items only \\
I10 & In the last 12 months, was there a time when, because of a lack of money or other resources, you went without eating for a whole day? & 1 Yes, 0 No & \texttt{fies\_whole\_day}\newline \textit{Skip: May be left blank}\newline Source: FAO Food Insecurity Experience Scale (FIES), SDG indicator 2.1.2; behavioural items only \\
\bottomrule
\end{longtable}
}
\subsection*{Module J: Ropeway}
\begin{quote}\small\textit{Read aloud: Now a few questions about the ropeway that has been proposed for this route. It is only a proposal — nothing has been built. There is no right answer and nothing you say here goes to anyone but us.}\end{quote}
{\scriptsize
\begin{longtable}{p{0.8cm}p{6.4cm}p{3.6cm}p{5.8cm}}
\toprule
\textbf{No.} & \textbf{Question or instruction} & \textbf{Answers} & \textbf{Variable, skip rule, source} \\
\midrule
\endfirsthead
\toprule
\textbf{No.} & \textbf{Question or instruction} & \textbf{Answers} & \textbf{Variable, skip rule, source} \\
\midrule
\endhead
J1 & A ropeway is proposed for this route. Are you in favour, neutral, or against? & 1 Support; 2 Neutral; 3 Oppose; 98 Prefer not to say & \texttt{ropeway\_stance}\newline Source: Project design \\
J2 & If the ropeway is built, what do you think will happen to your own work? More work for you; less work; about the same; or don't know? & 1 More work for me; 2 Less work for me; 3 About the same; 97 Don't know & \texttt{ropeway\_expect\_work}\newline Source: Pelz et al. 2024, Energy Policy 186:113973 (concept: perceived effect on livelihoods; wording ours) \\
J3 & Who do you think will get most of the jobs the ropeway creates? People from these villages; people from outside the area; nobody will get jobs; or don't know? & 1 People from these villages; 2 People from outside the area; 3 Nobody will get jobs; 97 Don't know & \texttt{ropeway\_expect\_jobs}\newline Source: Pelz et al. 2024, Energy Policy 186:113973 (concept: local vs outside benefit; wording ours) \\
J4 & Do you think the ropeway will improve the livelihoods of people here? Yes; no; or don't know? & 0 No; 1 Yes; 97 Don't know & \texttt{ropeway\_trust}\newline Source: Pelz et al. 2024, Energy Policy 186:113973 (Fig. 4B, trust in the company; wording ours) \\
J5 & Is your main work carrying, transporting or guiding people or goods on foot along the pilgrimage trek? & 1 Yes, 0 No & \texttt{trek\_dependent}\newline Source: Project design (see tasks\_module W4: exposure must be asked directly) \\
\bottomrule
\end{longtable}
}
\subsection*{Module K: Job quality (Apablaza et al. 2026 questions)}
\begin{quote}\small\textit{Read aloud: Now a set of questions about the conditions of your work — hours, pay, contract and so on. These come from a standard list used in many countries, so one or two may not fit your situation. Say so and we will move on.}\end{quote}
{\scriptsize
\begin{longtable}{p{0.8cm}p{6.4cm}p{3.6cm}p{5.8cm}}
\toprule
\textbf{No.} & \textbf{Question or instruction} & \textbf{Answers} & \textbf{Variable, skip rule, source} \\
\midrule
\endfirsthead
\toprule
\textbf{No.} & \textbf{Question or instruction} & \textbf{Answers} & \textbf{Variable, skip rule, source} \\
\midrule
\endhead
K1 & Do you have a signed contract? Yes, signed; yes but not yet signed; no contract. & 1 Yes, signed; 2 Yes, but not yet signed; 3 No contract & \texttt{contract\_status}\newline \textit{Skip: Ask only if employment\_type is 3 or 4 (wage workers) — self-employed skip to the next question}\newline Source: Apablaza et al. 2026, Appendix 2 Q11 \\
K2 & Is your workplace or business registered, for example with a taxpayer or GST number, a shop or trade licence, the Yatra registration, or a union? & 1 Yes, 0 No & \texttt{workplace\_registered}\newline Source: Apablaza et al. 2026, Appendix 2 Q12 (widened to local registrations) \\
K3 & Do you contribute to any pension system? Yes, the employer deducts it; yes, voluntarily; no. & 1 Yes, employer deducts it; 2 Yes, contributes voluntarily; 3 No & \texttt{pension\_contrib}\newline Source: Apablaza et al. 2026, Appendix 2 Q16 \\
K4 & Do you have health insurance through your work? Yes; only private or other insurance; none; don't know. & 1 Yes, through work; 2 Only private or other insurance; 3 None; 4 Don't know & \texttt{work\_health\_ins}\newline Source: Apablaza et al. 2026, Appendix 2 Q17 \\
K5 & Do you have the right to paid holiday, sick or maternity leave? & 0 No; 1 Yes; 97 Don't know & \texttt{leave\_rights}\newline Source: Apablaza et al. 2026, Appendix 2 Q18 \\
K6 & Have you ever been physically injured at your workplace? & 0 No; 1 Yes; 97 Don't know & \texttt{injured\_ever}\newline Source: Apablaza et al. 2026, Appendix 2 Q19 \\
K7 & In the last 12 months, was anyone physically injured at your workplace because of work? & 0 No; 1 Yes; 97 Don't know & \texttt{workplace\_injury\_12m}\newline Source: Apablaza et al. 2026, Appendix 2 Q20 (split into injury and death) \\
K8 & And in the last 12 months, did anyone die at your workplace because of work? & 0 No; 1 Yes; 97 Don't know & \texttt{workplace\_death\_12m}\newline Source: Apablaza et al. 2026, Appendix 2 Q20 (split into injury and death) \\
K9 & Would you like to work more hours than you do? & 0 No; 1 Yes; 97 Don't know & \texttt{wants\_more\_work}\newline Source: Apablaza et al. 2026, Appendix 2 Q21 \\
K10 & On a working day, how many MORE hours would you like to work? & Whole number & \texttt{more\_hours\_day}\newline \textit{Skip: Ask only if wants\_more\_work = 1}\newline Source: Apablaza et al. 2026, Appendix 2 Q22 \\
K11 & Across the months when you had no paid work, about how many MONTHS in total were you looking for work? & Whole number & \texttt{months\_looked\_for\_work}\newline \textit{Skip: Ask only if months\_no\_paid\_work > 0}\newline Source: Apablaza et al. 2026, Appendix 2 Q23 (duration in weeks, tied to the calendar instead of a single point in time) \\
K12 & Was your current work the first paid job you ever had? & 0 No; 1 Yes; 97 Don't know & \texttt{first\_job\_ever}\newline Source: Apablaza et al. 2026, Appendix 2 Q24 \\
K13 & Have you completed any college, diploma or university course? & 0 No; 1 Yes; 97 Don't know & \texttt{higher\_ed}\newline Source: Apablaza et al. 2026, Appendix 2 employment table, occupational status (1/8) \\
K14 & Are you enrolled in any pension, insurance or social-security scheme? For example, EPFO or PM-SYM pension, or Ayushman Bharat health cover. & 0 No; 1 Yes; 97 Don't know & \texttt{social\_security\_any}\newline Source: Apablaza et al. 2026, Appendix 2 employment table, employment security (1/8) \\
K15 & For more than half of your working day, do you work at very high speed? & 0 No; 1 Yes; 97 Don't know & \texttt{intens\_speed}\newline Source: Apablaza et al. 2026, Appendix 2 employment table, work intensity (1/16) \\
K16 & For more than half of your working day, do you work to tight deadlines? & 0 No; 1 Yes; 97 Don't know & \texttt{intens\_deadline}\newline Source: Apablaza et al. 2026, Appendix 2 employment table, work intensity (1/16) \\
K17 & For more than half of your working day, do you not have enough time to finish your tasks? & 0 No; 1 Yes; 97 Don't know & \texttt{intens\_time}\newline Source: Apablaza et al. 2026, Appendix 2 employment table, work intensity (1/16) \\
K18 & For more than half of your working day, do you work in a tiring or painful position? & 0 No; 1 Yes; 97 Don't know & \texttt{posture\_position}\newline Source: Apablaza et al. 2026, Appendix 2 employment table, posture risk (1/16) \\
K19 & For more than half of your working day, do you carry or move heavy loads? & 0 No; 1 Yes; 97 Don't know & \texttt{posture\_loads}\newline Source: Apablaza et al. 2026, Appendix 2 employment table, posture risk (1/16) \\
K20 & For more than half of your working day, do you make the same movements again and again? & 0 No; 1 Yes; 97 Don't know & \texttt{posture\_repetitive}\newline Source: Apablaza et al. 2026, Appendix 2 employment table, posture risk (1/16) \\
K21 & For more than half of your working day, are you exposed to loud noise? & 0 No; 1 Yes; 97 Don't know & \texttt{phys\_noise}\newline Source: Apablaza et al. 2026, Appendix 2 employment table, physical risk (1/16) \\
K22 & For more than half of your working day, are you exposed to extreme heat or cold? & 0 No; 1 Yes; 97 Don't know & \texttt{phys\_temperature}\newline Source: Apablaza et al. 2026, Appendix 2 employment table, physical risk (1/16) \\
\bottomrule
\end{longtable}
}
\subsection*{Module L: Tasks and skills (short block)}
\begin{quote}\small\textit{Read aloud: Now I will read out some kinds of work-tasks and ask whether you do them. There is no right or wrong answer — we are trying to understand what skills the work here actually uses.}\end{quote}
{\scriptsize
\begin{longtable}{p{0.8cm}p{6.4cm}p{3.6cm}p{5.8cm}}
\toprule
\textbf{No.} & \textbf{Question or instruction} & \textbf{Answers} & \textbf{Variable, skip rule, source} \\
\midrule
\endfirsthead
\toprule
\textbf{No.} & \textbf{Question or instruction} & \textbf{Answers} & \textbf{Variable, skip rule, source} \\
\midrule
\endhead
L1 & Do you load, unload, stack or count goods?  & 1 Yes, regularly in my main Yatra work; 2 Not in my main work, but done before elsewhere; 3 Never done & \texttt{tk\_load}\newline Source: Gathmann and Schonberg 2010; NCO 9333.0100 loader \\
L2 & Do you drive a motor vehicle?  & 1 Yes, regularly in my main Yatra work; 2 Not in my main work, but done before elsewhere; 3 Never done & \texttt{tk\_drive}\newline Source: Spitz-Oener 2006 / STEP; NCO 8322.0100 \\
L3 & Do you operate an engine or machine?  & 1 Yes, regularly in my main Yatra work; 2 Not in my main work, but done before elsewhere; 3 Never done & \texttt{tk\_engine}\newline Source: Gathmann and Schonberg 2010; NCO 8343.1700 ropeway operator \\
L4 & Do you do electrical or wiring work?  & 1 Yes, regularly in my main Yatra work; 2 Not in my main work, but done before elsewhere; 3 Never done & \texttt{tk\_electric}\newline Source: Spitz-Oener 2006; NCO 7411.0301 wireman \\
L5 & Do you check equipment or the route for safety?  & 1 Yes, regularly in my main Yatra work; 2 Not in my main work, but done before elsewhere; 3 Never done & \texttt{tk\_safety}\newline Source: Gathmann and Schonberg 2010; NCO 8343.1700 \\
L6 & Do you sell goods or services?  & 1 Yes, regularly in my main Yatra work; 2 Not in my main work, but done before elsewhere; 3 Never done & \texttt{tk\_sell}\newline Source: Gathmann and Schonberg 2010; NCO 5223.0200 \\
L7 & Do you handle cash and payments?  & 1 Yes, regularly in my main Yatra work; 2 Not in my main work, but done before elsewhere; 3 Never done & \texttt{tk\_cash}\newline Source: Autor and Handel 2013; NCO 5131.0401 \\
L8 & Do you cook or prepare food or drink?  & 1 Yes, regularly in my main Yatra work; 2 Not in my main work, but done before elsewhere; 3 Never done & \texttt{tk\_cook}\newline Source: NCO 5120.0300 cook \\
L9 & Do you serve guests or customers?  & 1 Yes, regularly in my main Yatra work; 2 Not in my main work, but done before elsewhere; 3 Never done & \texttt{tk\_serve}\newline Source: NCO 5131.0401 waiter; Pal et al. 2026 \\
L10 & Do you clean rooms or public areas?  & 1 Yes, regularly in my main Yatra work; 2 Not in my main work, but done before elsewhere; 3 Never done & \texttt{tk\_clean}\newline Source: NCO 5151.0202 room attendant; Autor-Levy-Murnane 2003 (janitorial) \\
L11 & Do you guide or explain things to visitors?  & 1 Yes, regularly in my main Yatra work; 2 Not in my main work, but done before elsewhere; 3 Never done & \texttt{tk\_guide}\newline Source: Gathmann and Schonberg 2010; NCO 5113.0200 \\
L12 & Do you coordinate work by phone, radio or signals?  & 1 Yes, regularly in my main Yatra work; 2 Not in my main work, but done before elsewhere; 3 Never done & \texttt{tk\_coord}\newline Source: NCO 8343.1700 ropeway operator \\
L13 & Have you ever driven a motorcycle or scooter, with or without a licence? & 1 Yes, 0 No & \texttt{drove\_twowheeler}\newline \textit{Skip: Ask only if tk\_drive is 1 or 2}\newline Source: Level ladder without licences: people drive without holding a licence \\
L14 & ... a car or jeep? & 1 Yes, 0 No & \texttt{drove\_car}\newline \textit{Skip: Ask only if tk\_drive is 1 or 2}\newline Source: Level ladder without licences \\
L15 & ... a truck or bus? & 1 Yes, 0 No & \texttt{drove\_heavy}\newline \textit{Skip: Ask only if tk\_drive is 1 or 2}\newline Source: Level ladder without licences \\
L16 & In your work, do you read anything (notes, rate lists, tickets, messages)? & 1 Yes, 0 No & \texttt{read\_at\_work}\newline Source: STEP module 6A \\
L17 & In your work, do you work out prices or costs? & 1 Yes, 0 No & \texttt{calc\_at\_work}\newline Source: STEP module 6A \\
\bottomrule
\end{longtable}
}
\subsection*{Module M: Access to common resources (short block)}
\begin{quote}\small\textit{Read aloud: Last, a few questions about the forest and common land around where you work — firewood, fodder, grazing and water. Some of this is not allowed in some places, and that is a perfectly good answer: we only want to know what is possible and what is not. Nothing you say goes to the Forest Department or to anyone else, and you can leave any of these out.}\end{quote}
{\scriptsize
\begin{longtable}{p{0.8cm}p{6.4cm}p{3.6cm}p{5.8cm}}
\toprule
\textbf{No.} & \textbf{Question or instruction} & \textbf{Answers} & \textbf{Variable, skip rule, source} \\
\midrule
\endfirsthead
\toprule
\textbf{No.} & \textbf{Question or instruction} & \textbf{Answers} & \textbf{Variable, skip rule, source} \\
\midrule
\endhead
M1 & Where you work, is it allowed to take firewood, fodder or forest produce from the forest or common land? Ask what they understand the rule to be; do not correct them. & 1 Yes, it is allowed; 0 No, it is not allowed; 2 Allowed only with a permit or a fee; 97 Don't know & \texttt{cpr\_permitted}\newline \textit{Skip: May be left blank}\newline Source: Project design (de jure access, 2026-10-05); Ostrom 1990 on the distinction between rule in form and rule in use \\
M2 & In this Yatra season or last Yatra season, did you or anyone staying with you here collect firewood or dead wood from common land or forest in the Yatra region or workplace region? & 0 No; 1 Yes; 2 No, it is not allowed here; 97 Don't know & \texttt{cpr\_firewood}\newline \textit{Skip: May be left blank}\newline Source: Krantz 2001, natural capital: resource stocks from which resource flows are derived (definitions list); item wording is ours \\
M3 & In this Yatra season or last Yatra season, did you or anyone staying with you here cut grass or fodder from common land in the Yatra region or workplace region? & 0 No; 1 Yes; 2 No, it is not allowed here; 97 Don't know & \texttt{cpr\_fodder}\newline \textit{Skip: May be left blank}\newline Source: Krantz 2001, natural capital: resource stocks from which resource flows are derived; item wording is ours \\
M4 & In this Yatra season or last Yatra season, did you or anyone staying with you here collect herbs, mushrooms or wild fruit from the forest in the Yatra region or workplace region? & 0 No; 1 Yes; 2 No, it is not allowed here; 97 Don't know & \texttt{cpr\_forest\_produce}\newline \textit{Skip: May be left blank}\newline Source: Krantz 2001, natural capital (resource flows from stocks); item wording is ours \\
M5 & In this Yatra season or last Yatra season, did any animal of yours or of anyone staying with you here graze on common pasture in the Yatra region or workplace region? & 0 No; 1 Yes; 2 No, it is not allowed here; 97 Don't know & \texttt{cpr\_grazing}\newline \textit{Skip: May be left blank}\newline Source: Krantz 2001, natural capital: resource stocks from which resource flows are derived; item wording is ours \\
M6 & In this Yatra season or last Yatra season, has anyone you or someone staying with you here been stopped from using a common forest, pasture or path, or been asked to pay to use it, in the Yatra region or workplace region? & 0 No; 1 Yes; 97 Don't know & \texttt{cpr\_restricted}\newline \textit{Skip: May be left blank}\newline Source: Krantz 2001, access ('the opportunity in practice to use a resource') and claims ('demands and appeals'); item wording is ours \\
M7 & In this Yatra season or last Yatra season, has a common place you used to collect from or graze on been closed, because of a landslide, flood, road or bridge damage, or construction work, in the Yatra region or workplace region? & 0 No; 1 Yes; 97 Don't know & \texttt{cpr\_lost\_access}\newline \textit{Skip: May be left blank}\newline Source: Krantz 2001, access (as above); the ropeway link is the project's own \\
M8 & On your route or near your worksite, can you use a common water source (spring, tap or stream) without paying? & 1 Yes, free to use; 2 A source exists but we must pay; 3 No source nearby; 97 Don't know & \texttt{cpr\_water\_route}\newline \textit{Skip: May be left blank}\newline Source: Krantz 2001, natural capital (water) and access ('the opportunity in practice to use a resource'); item wording is ours \\
M9 & In the Yatra region or workplace region, who decides who may use the common land and forest? Ask and code what they name; do not read the list. & 1 Village head or panchayat; 2 Forest department; 3 Local committee or community group; 4 Nobody decides; it is open to all; 97 Don't know & \texttt{cpr\_governance}\newline \textit{Skip: May be left blank}\newline Source: NO SOURCE ON DISK YET: Krantz does not address governance of common land. To be sourced from Agrawal 2001 or Ostrom 1990 once read on disk \\
\bottomrule
\end{longtable}
}

\section*{4. Variables built afterwards (not asked)}
All of these are computed in Stata (\texttt{do/02\_build\_final\_dataset.do}) from the asked answers, so the analysis starts from raw answers and every step can be checked.
{\scriptsize
\begin{longtable}{p{3.3cm}p{7.6cm}p{5.7cm}}
\toprule
\textbf{Variable} & \textbf{Meaning} & \textbf{How it is made} \\
\midrule
\endfirsthead
\toprule
\textbf{Variable} & \textbf{Meaning} & \textbf{How it is made} \\
\midrule
\endhead
\texttt{education\_years} & The direct year count (years\_schooling) when the respondent knew it; otherwise the midpoint of the fallback bracket (education\_level\_cat: 0/2/7/11 for no-formal/some-primary/primary-complete/secondary-plus); missing if the respondent preferred not to say. & years\_schooling if knows\_years\_schooling==1, else bracket midpoint from education\_level\_cat \\
\texttt{hoh\_female} & Derived from the gendered relationship categories: husband/father/son/brother/other-male imply a male head, wife/mother/daughter/sister/other-female a female head, and when the respondent IS the head it is their own sex. Replaces a separate male/female question. & from hoh\_relation; = female if hoh\_relation==1 \\
\texttt{months\_here} & Derived from the absence spell: 12 minus the months between left\_here\_month and returned\_here\_month, or 12 for a respondent who never leaves. This is the measured length of the respondent OWN season, in place of a constant applied to everyone, and it is what the mid-month Yatra-start problem resolves to — the boundary is the month they reported, not one we chose. Built straight from the spell as of 2026-10-01; it used to be 12 minus months\_home\_base minus months\_third\_place, which routed the season weight — the weight on every consumption and income figure in the study — through the months\_away\_for\_work answer. It no longer depends on it. & 12 - months\_away\_total \\
\texttt{closure\_labour\_migrant} & 1 if the respondent reported working away during the closure. Type B in the closure-regime typology and the mobility variable that actually carries information in this population, in place of a migrant dummy built on district boundaries — which the pilot shows would classify 11 of 20 seasonal movers as non-movers. & worked\_away\_in\_closure==1 \\
\texttt{cons\_pc\_denom\_season} & n\_here\_season for a split household, hhsize for everyone else. This is the fix for a real error in the poverty headcount: per-capita consumption divided the Yatra-season figure — which covers only "you and anyone staying with you here" — by the FULL household size, so a man supporting himself here for six months while a family of five lived at the home place was recorded at a fraction of his true per-capita consumption and counted as poor by arithmetic. The off-season half keeps hhsize, because by then the household is reunited. & hhsize \\
\texttt{stays\_all\_year} & 1 if the respondent does not go to the home place at closure. The case the two-season design of this instrument does NOT fit: for these respondents the Yatra-season and off-season questions describe the same place, and the pair should be checked for a suspiciously high identical-answer rate at the pilot. & resp\_returns\_at\_closure==0 \\
\texttt{split\_household} & 1 if the rest of the household lives at the home place year-round. The case Module E Yatra-season wording was written for, and the case whose per-capita consumption was being computed wrongly until n\_here\_season existed. & hh\_at\_home\_place==1 \\
\texttt{credit\_institutional} & 1 if the household borrowed in the last 12 months from a bank, cooperative, RRB, microfinance institution or SHG (credit\_source 1-4); 0 otherwise, INCLUDING households that did not borrow at all. Defined for every respondent on purpose: credit\_source itself is now gated behind took\_loan\_12m, so using it directly as a VEP covariate would drop every non-borrower from the regression. & took\_loan\_12m==1 \& inrange(credit\_source,1,4) \\
\texttt{credit\_informal} & 1 if the household borrowed from a moneylender, relative, friend, employer or contractor (credit\_source 5-8); 0 otherwise, including non-borrowers. Separated from institutional credit because the two have opposite signs for vulnerability. & took\_loan\_12m==1 \& inrange(credit\_source,5,8) \\
\texttt{migrant} & 1 if origin is not Local (same district). & origin > 1 \\
\texttt{health\_access\_tier} & Remoteness of the interview site along the route, derived from the interview GPS rather than from an enumerator-coded cluster: the route runs from the road-head and its hospital up to the shrine, so position along it is what determines how reachable care is. Latitude bands (research team's route crosswalk): below 30.58 Good, 30.58 to 30.66 Moderate, 30.66 and above Poor. Bands, not coordinates, enter the analysis, so the exported data still carries no named site. & from gps\_lat: 3 if <30.58; 2 if 30.58-30.66; 1 if >=30.66 \\
\texttt{health\_access\_deprived} & 1 if the site tier is Poor, or Moderate and nobody in the household has health cover. health\_insurance\_covered was REMOVED: it asked whether "your household" is covered, which for a split migrant household — the worker here, the family elsewhere — does not name a group the respondent can answer for. n\_health\_insured, a count of covered members, replaces it and is strictly more informative. This variable is also no longer part of the MPI: the health dimension now uses NITI's own mortality and maternal indicators, so this is kept only as a VEP exposure covariate. & tier==1 | (tier==2 \& n\_health\_insured==0) \\
\texttt{yatra\_months} & Count of calendar months coded 1, Yatra work proper. Preparation months (status 9) are NOT counted: this is the denominator for earnings per Yatra month, and a month spent buying stock and repairing tack is not a month that produced the income being divided. & months\_worked\_yatra, asked directly \\
\texttt{yatra\_months\_liv} & yatra\_months + yatra\_prep\_months, i.e. exactly what yatra\_months counted before the 2026-10-03 split. Every season-weighted annual figure below uses THIS, not yatra\_months, so those figures are numerically unchanged by the split: the respondent is tied to the route during preparation as much as during the season, and it is that, not whether he was earning, that decides which spending and coping regime he is in. & yatra\_months + yatra\_prep\_months \\
\texttt{offseason\_months\_worked} & Months coded 2 to 7 (any paid activity other than Yatra work). & 12 - months\_worked\_yatra - months\_no\_paid\_work, floored at 0 \\
\texttt{months\_no\_work} & Months coded 8. & months\_no\_paid\_work, asked directly \\
\texttt{offseason\_primary} & The activity that fills most of the non-Yatra months (ties go to the lower code; all months idle gives 8). & offseason\_activity, asked directly \\
\texttt{yatra\_income} & Calendar path: sum of monthly earnings in the months coded Yatra work. Fallback path: the annual total times the respondent's own Yatra share. & income\_annual\_total * pct\_income\_yatra/100 \\
\texttt{non\_yatra\_income} & Calendar path: sum of monthly earnings in the other months. Fallback path: the annual total times the remaining share. & income\_annual\_total * (100-pct\_income\_yatra)/100 \\
\texttt{income\_pm\_yatra\_eq} & Yatra earnings spread evenly across the months the calendar codes as Yatra work. Defined for fallback respondents only; for calendar respondents the actual monthly figures are used instead. & yatra\_income / yatra\_months if income\_from\_fallback==1 \\
\texttt{income\_pm\_other\_eq} & Non-Yatra earnings spread evenly across the remaining months. Defined for fallback respondents only. & non\_yatra\_income / (12 - yatra\_months\_liv) if income\_from\_fallback==1 \\
\texttt{remittance\_outward\_annual} & Yatra-season months at the Yatra-season usual amount, other months at the off-season usual amount. & yatra\_months\_liv*remit\_out\_yatra\_pm + (12-yatra\_months\_liv)*remit\_out\_offseason\_pm \\
\texttt{remittance\_inward\_annual} & Yatra-season months at the Yatra-season usual amount, other months at the off-season usual amount. & yatra\_months\_liv*remit\_in\_yatra\_pm + (12-yatra\_months\_liv)*remit\_in\_offseason\_pm \\
\texttt{total\_annual\_income} & Work income plus inward remittances. Outward remittances are not subtracted. & yatra\_income + non\_yatra\_income + remittance\_inward\_annual \\
\texttt{yatra\_income\_share} & Yatra income divided by total annual income (an Exposure covariate in the VEP model). & yatra\_income / total\_annual\_income \\
\texttt{income\_seasonality\_cv} & Calendar path: coefficient of variation across the twelve reported monthly earnings (n-1 divisor). Fallback path: the same CV computed on a two-level series — income\_pm\_yatra\_eq in each Yatra month, income\_pm\_other\_eq in each other month — so the variable is never missing and no respondent silently drops out of the FGLS for want of it. IMPORTANT: the fallback series has no within-season variation by construction, so its CV captures only the between-season swing and is a LOWER BOUND on true seasonality. Always run the VEP arms with income\_from\_fallback interacted or split; do not treat the two paths as the same measurement. & calendar: rowsd(income\_m1..12)/rowmean(income\_m1..12); fallback: CV of the two-level series built from income\_pm\_yatra\_eq and income\_pm\_other\_eq over yatra\_months and 12-yatra\_months \\
\texttt{cons\_staples\_pm} & Staples spending, weighted by the length of the Yatra season and the rest of the year. & (yatra\_months\_liv*cons\_staples\_yatra\_pm + (12-yatra\_months\_liv)*cons\_staples\_offseason\_pm)/12 \\
\texttt{cons\_perishables\_pm} & Perishables spending, weighted by the length of the Yatra season and the rest of the year. & (yatra\_months\_liv*cons\_perishables\_yatra\_pm + (12-yatra\_months\_liv)*cons\_perishables\_offseason\_pm)/12 \\
\texttt{cons\_food\_own\_pm} & Home-grown/gifted food spending, weighted by the length of the Yatra season and the rest of the year. & (yatra\_months\_liv*cons\_food\_own\_yatra\_pm + (12-yatra\_months\_liv)*cons\_food\_own\_offseason\_pm)/12 \\
\texttt{cons\_food\_out\_pm} & Food eaten outside spending, weighted by the length of the Yatra season and the rest of the year. & (yatra\_months\_liv*cons\_food\_out\_yatra\_pm + (12-yatra\_months\_liv)*cons\_food\_out\_offseason\_pm)/12 \\
\texttt{cons\_fuel\_pm} & Fuel and light spending, weighted by the length of the Yatra season and the rest of the year. & (yatra\_months\_liv*cons\_fuel\_yatra\_pm + (12-yatra\_months\_liv)*cons\_fuel\_offseason\_pm)/12 \\
\texttt{cons\_routine\_misc\_pm} & Toiletries/consumables spending, weighted by the length of the Yatra season and the rest of the year. & (yatra\_months\_liv*cons\_routine\_misc\_yatra\_pm + (12-yatra\_months\_liv)*cons\_routine\_misc\_offseason\_pm)/12 \\
\texttt{cons\_transport\_comm\_pm} & Transport/communication spending, weighted by the length of the Yatra season and the rest of the year. & (yatra\_months\_liv*cons\_transport\_comm\_yatra\_pm + (12-yatra\_months\_liv)*cons\_transport\_comm\_offseason\_pm)/12 \\
\texttt{cons\_rent\_pm} & Rent spending, weighted by the length of the Yatra season and the rest of the year. & (yatra\_months\_liv*cons\_rent\_yatra\_pm + (12-yatra\_months\_liv)*cons\_rent\_offseason\_pm)/12 \\
\texttt{cons\_med\_nonhosp\_pm} & Medical, non-hospital spending, weighted by the length of the Yatra season and the rest of the year. & (yatra\_months\_liv*cons\_med\_nonhosp\_yatra\_pm + (12-yatra\_months\_liv)*cons\_med\_nonhosp\_offseason\_pm)/12 \\
\texttt{cons\_packaged\_food\_pm} & Packaged snacks and drinks spending, weighted by the length of the Yatra season and the rest of the year. & (yatra\_months\_liv*cons\_packaged\_food\_yatra\_pm + (12-yatra\_months\_liv)*cons\_packaged\_food\_offseason\_pm)/12 \\
\texttt{cons\_pan\_tobacco\_pm} & Pan, tobacco and intoxicants spending, weighted by the length of the Yatra season and the rest of the year. & (yatra\_months\_liv*cons\_pan\_tobacco\_yatra\_pm + (12-yatra\_months\_liv)*cons\_pan\_tobacco\_offseason\_pm)/12 \\
\texttt{cons\_food\_pm} & Staples + perishables + home-grown/gifted + eaten outside + packaged snacks and drinks, each already weighted by season. Packaged food is inside HCES's own food block (Section 7), so it belongs here; pan, tobacco and intoxicants are a separate MPCE category in HCES and are added to total\_cons\_pm instead. & cons\_staples\_pm + cons\_perishables\_pm + cons\_food\_own\_pm + cons\_food\_out\_pm + cons\_packaged\_food\_pm \\
\texttt{cons\_clothing\_12m\_pm} & Annual figure divided by 12, rounded. & round(cons\_clothing\_12m/12) \\
\texttt{cons\_education\_12m\_pm} & Annual figure divided by 12, rounded. & round(cons\_education\_12m/12) \\
\texttt{cons\_medical\_12m} & Non-hospitalisation (annual-average month x 12) + hospitalisation (12-month). & cons\_med\_nonhosp\_pm*12 + cons\_medical\_hosp\_12m \\
\texttt{cons\_medical\_12m\_pm} & Non-hospitalisation (already an annual-average month) + hospitalisation divided by 12, rounded. & cons\_med\_nonhosp\_pm + round(cons\_medical\_hosp\_12m/12) \\
\texttt{cons\_durables\_12m\_pm} & Annual figure divided by 12, rounded. & round(cons\_durables\_12m/12) \\
\texttt{total\_cons\_pm} & Food (annual-average month above) + fuel + toiletries + transport/communication + rent (each annual-average) + the four monthly-equivalent annual items. This is a true annual-average monthly figure, not a single-season snapshot. & food\_pm (now incl. packaged) + pan\_tobacco\_pm + fuel\_pm + routine\_misc\_pm + transport\_comm\_pm + rent\_pm + clothing\_pm + education\_pm + medical\_pm + durables\_pm \\
\texttt{cons\_pc\_pm} & Household consumption per month divided by household size, rounded. This is the VEP outcome (its log). & round(total\_cons\_pm / hhsize) \\
\texttt{cons\_pc\_pm\_narrow} & Only the non-lumpy items (food incl. packaged, pan/tobacco, fuel, toiletries, transport/communication, rent), each already an annual-average month, per person. Leaves out the annual items, as in Dercon and Krishnan 2000. For robustness of the variance model; not comparable with the poverty line. & (food\_pm + pan\_tobacco\_pm + fuel\_pm + routine\_misc\_pm + transport\_comm\_pm + rent\_pm) / hhsize \\
\texttt{n\_children\_u15} & The two asked bands added: children under 6 plus children aged 6 to 14. Asked directly until 2026-10-01, alongside the 6-to-14 count that is a subset of it — which made the respondent subtract one from the other and let the two contradict. Two disjoint bands partition it exactly, so this is now arithmetic and the contradiction cannot arise. & n\_children\_u6 + n\_children\_6\_14 \\
\texttt{cons\_pc\_ae\_pm} & Household consumption divided by (adults + 0.5 x children under 15), where adults = household size minus children under 15. One simple scale; the choice of scale is a sensitivity check. & total\_cons\_pm / (hhsize - n\_children\_u15 + 0.5*n\_children\_u15) \\
\texttt{poverty\_line} & Sethu, Surya and Ruthu (2024) construct TWO lines for 2022-23 by the Rangarajan method: Rs 2,515 per capita per month rural and Rs 3,639 urban. Which applies is decided by home\_rural\_urban — the respondent's USUAL place of residence, not the Yatra worksite. Until this variable existed the build applied the rural line to everyone, which measured urban-resident workers against a line 31 percent too low and understated their poverty. & 2515 if home\_rural\_urban==1, else 3639 \\
\texttt{poor} & 1 if per-capita monthly consumption is below the rural or urban line that applies to this respondent. & cons\_pc\_pm < poverty\_line \\
\texttt{poor\_sensitivity\_cpi} & 1 if below the Rangarajan and Dev 2024 CPI-adjusted line; robustness only. & cons\_pc\_pm < 1850 \\
\texttt{durables\_count} & Legacy six-item count kept only so earlier outputs stay reproducible. DO NOT use for the MPI: the asset indicator is mpi\_asset\_deprived, which applies NITI's actual rule to its actual eight-item list. & owns\_tv + owns\_radio + owns\_bicycle + owns\_motorcycle + owns\_car + owns\_fridge \\
\texttt{productive\_assets\_count} & Sum of the six productive-asset binaries (livestock, pony/mule, shop/stall, work vehicle, work equipment). & owns\_cow\_buffalo + owns\_goat\_sheep + owns\_pony\_mule + owns\_shop\_stall + owns\_work\_vehicle + owns\_work\_equipment \\
\texttt{water\_deprived} & NITI rule, BOTH limbs: deprived if the source is unimproved (codes 7-9), OR the source is improved but the round trip to fetch it takes over 30 minutes. The second limb could not be evaluated at all before water\_on\_premises and water\_fetch\_minutes existed. & inrange(drinking\_water,7,9) | (water\_on\_premises==0 \& water\_fetch\_minutes>30) \\
\texttt{cooking\_fuel\_deprived} & NITI rule: deprived if the household cooks with firewood, dung, crop residue/shrubs, charcoal or coal (codes 6-10). Kerosene (5) is NOT counted, because NITI's own list does not name it; the global MPI does count it, so this is a documented departure and the raw fuel code is retained so either rule can be applied later. & inlist(cooking\_fuel,6,7,8,9,10) \\
\texttt{mpi\_asset\_count} & Radio, TV, telephone, computer, animal cart, bicycle, motorbike, refrigerator — NITI's own list. Car or truck is NOT in this count; it is a separate limb of the indicator. & owns\_radio + owns\_tv + owns\_phone + owns\_computer + owns\_animal\_cart + owns\_bicycle + owns\_motorcycle + owns\_fridge \\
\texttt{mpi\_asset\_deprived} & NITI rule exactly: deprived if the household does not own MORE THAN ONE of the eight small assets AND does not own a car or truck. The previous version summed six assets into a plain count, which is not NITI's rule and is not comparable to any published National MPI figure. & mpi\_asset\_count <= 1 \& owns\_car == 0 \\
\texttt{shock\_1} & 1 if 'Serious illness or death of an earning member' was ticked. & distress\_event\_last365d contains 1 \\
\texttt{shock\_2} & 1 if 'Crop failure or livestock loss' was ticked. & distress\_event\_last365d contains 2 \\
\texttt{shock\_3} & 1 if 'Damage from a natural disaster (flood, landslide, fire)' was ticked. & distress\_event\_last365d contains 3 \\
\texttt{shock\_4} & 1 if 'Loss of business, goods or assets' was ticked. & distress\_event\_last365d contains 4 \\
\texttt{shock\_5} & 1 if 'Road or bridge blocked, route closed' was ticked. & distress\_event\_last365d contains 5 \\
\texttt{shock\_6} & 1 if 'Yatra stopped early or badly disrupted' was ticked. & distress\_event\_last365d contains 6 \\
\texttt{shock\_7} & 1 if 'Sharp fall in customers or prices' was ticked. & distress\_event\_last365d contains 7 \\
\texttt{shock\_8} & 1 if 'Lost the work or the work stopped' was ticked. & distress\_event\_last365d contains 8 \\
\texttt{no\_shock\_reported} & 1 if code 9 was ticked. Distinct from an empty list, which the form can no longer produce: the item is required and code 9 is how a household with no shock answers it. & distress\_event\_last365d contains 9 \\
\texttt{cope\_1} & 1 if 'Used savings' was ticked. & shock\_coping contains 1 \\
\texttt{cope\_2} & 1 if 'Borrowed money' was ticked. & shock\_coping contains 2 \\
\texttt{cope\_3} & 1 if 'Sold or pawned assets' was ticked. & shock\_coping contains 3 \\
\texttt{cope\_4} & 1 if 'Cut consumption' was ticked. & shock\_coping contains 4 \\
\texttt{cope\_5} & 1 if 'Help from relatives/community' was ticked. & shock\_coping contains 5 \\
\texttt{cope\_6} & 1 if 'Paid it out of normal earnings, nothing given up' was ticked. & shock\_coping contains 6 \\
\texttt{cope\_7} & 1 if 'Something else (say what)' was ticked. & shock\_coping contains 7 \\
\texttt{shock\_count} & Count of shock types ticked. The old single-answer item could only ever be 0 or 1. & sum of shock\_1..8 \\
\texttt{shock\_any} & 1 if any shock was ticked. & shock\_count > 0 \\
\texttt{shock\_covariate} & 1 if any of: natural disaster, road or bridge blocked, Yatra stopped or disrupted, sharp fall in customers or prices (codes 3, 5, 6, 7). These hit the whole route at once, so they are the shocks a ropeway — or a flood year — delivers to everyone simultaneously. Separating them from idiosyncratic shocks is what makes the covariate/idiosyncratic variance decomposition identifiable. & any of shock\_3 shock\_5 shock\_6 shock\_7 \\
\texttt{shock\_idiosyncratic} & 1 if any of: illness or death of an earning member, crop or livestock loss, loss of business or assets, lost the work (codes 1, 2, 4, 8). These hit one household at a time. & any of shock\_1 shock\_2 shock\_4 shock\_8 \\
\texttt{coped\_sold\_assets} & 1 if selling or pawning assets was among the responses. Kept as its own variable because the asset-smoothing test compares it directly against cutting consumption, and the two are no longer mutually exclusive. & cope\_3 \\
\texttt{coped\_cut\_consumption} & 1 if cutting consumption was among the responses. See coped\_sold\_assets — a household doing BOTH, or cutting consumption precisely in order to avoid selling, is the case the single-answer item could not represent. & cope\_4 \\
\texttt{mpi\_mortality\_dep} & NITI: a child or adolescent under 18 died in the household in the last five years. & child\_death\_5y == 1 \\
\texttt{mpi\_maternal\_dep} & NITI: a woman who gave birth in the last five years did not have at least four antenatal visits for the most recent birth, OR was not assisted by trained personnel at that birth. Households with no birth in the window are NOT deprived, per NITI's own treatment. & birth\_last\_5y==1 \& (anc\_4\_visits!=1 | !inlist(skilled\_birth\_attendant,1,2)) \\
\texttt{mpi\_schooling\_dep} & NITI: not even one household member aged 10 or older has completed six years of schooling. & any\_member\_6yr\_schooling == 0 \\
\texttt{mpi\_attendance\_dep} & NITI: any school-aged child is not attending school up to the age at which they would complete class 8. Not deprived where the household has no child in that range. & (n\_children\_6\_14 - n\_children\_in\_school) > 0 if n\_children\_6\_14 > 0 \\
\texttt{mpi\_housing\_dep} & NITI: the floor is natural material, OR the roof OR the wall is rudimentary. All three limbs are now collected; the wall limb was previously missing entirely. & floor\_material==1 | roof\_material==1 | wall\_material==1 \\
\texttt{mpi\_sanitation\_dep} & NITI: unimproved or no facility, OR improved but shared with other households. & inlist(toilet\_type,1,2,3) \\
\texttt{mpi\_electricity\_dep} & NITI: the household has no electricity. & electricity == 0 \\
\texttt{mpi\_bank\_dep} & NITI: no household member has a bank account or a post office account. & has\_bank\_account == 0 \\
\texttt{mpi\_score} & Weighted sum of the eleven collected indicators — mortality, maternal health, years of schooling, school attendance, cooking fuel, sanitation, drinking water, electricity, housing, assets and bank account. Nutrition (1/6) cannot be collected without anthropometry, so the remaining weights are scaled up proportionally to sum to 1. REPORT THIS AS AN ADAPTATION, never as the National MPI, and name the absent indicator. mpi\_score\_lyons is the companion that substitutes rather than omits. (The label said TEN until 2026-10-01; eleven of the twelve are collected, and the formula always summed eleven.) & (1/12*mortality + 1/12*maternal + 1/6*schooling + 1/6*attendance + 1/21*(fuel+sanitation+water+electricity+housing+assets+bank)) / (1 - 1/6) \\
\texttt{mpi\_poor} & Alkire-Foster identification at the standard k = 33.3 percent cutoff. & mpi\_score >= 1/3 \\
\texttt{shock\_earnings\_lost} & Weeks of work lost times the respondent's own weekly earnings in the season the shock fell in — Yatra-season weekly earnings if shock\_month is inside the measured season, off-season weekly earnings otherwise. The respondent is never asked to value his own forgone work: he reports the weeks, and the wage is the one this instrument already measured from him. & shock\_work\_lost\_weeks * (season-matched weekly earnings) \\
\texttt{shock\_loss\_total} & The two limbs added. This is what shock\_loss\_amount used to ask the respondent to add up in his head — including valuing his own lost work — and it replaces it everywhere, including in shock\_loss\_share. Report the limbs separately too: a shock that costs only time and a shock that costs only cash are different events for a household with no savings. & shock\_earnings\_lost + shock\_money\_spent \\
\texttt{owns\_tv} & Split out of assets\_owned; falls back to the asked column on records collected before 2026-10-05. & assets\_owned contains 1 \\
\texttt{owns\_radio} & Split out of assets\_owned; falls back to the asked column on records collected before 2026-10-05. & assets\_owned contains 2 \\
\texttt{owns\_bicycle} & Split out of assets\_owned; falls back to the asked column on records collected before 2026-10-05. & assets\_owned contains 3 \\
\texttt{owns\_motorcycle} & Split out of assets\_owned; falls back to the asked column on records collected before 2026-10-05. & assets\_owned contains 4 \\
\texttt{owns\_car} & Split out of assets\_owned; falls back to the asked column on records collected before 2026-10-05. & assets\_owned contains 5 \\
\texttt{owns\_phone} & Split out of assets\_owned; falls back to the asked column on records collected before 2026-10-05. & assets\_owned contains 6 \\
\texttt{owns\_computer} & Split out of assets\_owned; falls back to the asked column on records collected before 2026-10-05. & assets\_owned contains 7 \\
\texttt{owns\_animal\_cart} & Split out of assets\_owned; falls back to the asked column on records collected before 2026-10-05. & assets\_owned contains 8 \\
\texttt{owns\_fridge} & Split out of assets\_owned; falls back to the asked column on records collected before 2026-10-05. & assets\_owned contains 9 \\
\texttt{owns\_cow\_buffalo} & Split out of livestock\_owned; falls back to the asked column on records collected before 2026-10-05. & livestock\_owned contains 1 \\
\texttt{owns\_goat\_sheep} & Split out of livestock\_owned; falls back to the asked column on records collected before 2026-10-05. & livestock\_owned contains 2 \\
\texttt{owns\_pony\_mule} & Split out of livestock\_owned; falls back to the asked column on records collected before 2026-10-05. & livestock\_owned contains 3 \\
\texttt{fies\_raw} & Count of yes answers across the five FIES behavioural items. Unweighted, by design: FIES is an ordered severity scale, not a frequency-by-severity index like the rCSI it replaces, so applying weights to it would be a category error. Missing only where all five were left blank. & fies\_skipped\_meal + fies\_ate\_less + fies\_ran\_out + fies\_hungry + fies\_whole\_day \\
\texttt{fies\_mod\_sev} & 1 if two or more of the five behavioural items were answered yes. Replaces food\_coping\_deprived (rCSI > 20), whose threshold was Lyons et al.'s, calibrated on Lebanese refugee households answering the original 7-day rCSI -- a cutoff on a quantity this instrument was not measuring. Two or more behavioural items is the conventional moderate-or-severe line. & fies\_raw >= 2 \\
\texttt{mpi\_nutrition\_proxy\_dep} & fies\_mod\_sev, standing in for NITI's Nutrition indicator in mpi\_score\_lyons only. NOT a measure of undernourishment and must never be labelled as one: it is moderate-or-severe food insecurity on the FAO FIES behavioural items, occupying the 1/6 weight NITI gives to a measurement this instrument cannot take (anthropometry). Rewired from food\_coping\_deprived (rCSI > 20) on 2026-10-05 with the rest of the food-security block. FIES is the better occupant of the slot on its own terms -- it measures experience of insufficient food, which is one step from undernourishment, where the rCSI measured coping acts, which is two -- but it is still a substitution and must be reported as one. & fies\_mod\_sev == 1 \\
\texttt{mpi\_score\_lyons} & All twelve NITI weights filled, with the food-coping indicator in the Nutrition slot at 1/6 and no reweighting. Report alongside mpi\_score, never instead of it, and state the substitution. & mpi\_score with (1/6)*mpi\_nutrition\_proxy\_dep in place of the reweighting \\
\texttt{mpi\_poor\_lyons} & Alkire-Foster identification at k = 33.3 percent on mpi\_score\_lyons. The gap between this headcount and mpi\_poor is the sensitivity of the result to the nutrition substitution. & mpi\_score\_lyons >= 1/3 \\
\texttt{hours\_week\_yatra} & Hours per day times days per week. & hours\_day\_yatra * days\_week\_yatra \\
\texttt{hours\_week\_offseason} & Hours per day times days per week, off-season. Missing for respondents with no off-season paid work at all. & hours\_day\_offseason * days\_week\_offseason \\
\texttt{hours\_week\_annual} & Yatra-season and off-season weekly hours, weighted by how many months of each the calendar records. Averaged over WORKED months only, so idle months do not drag the figure to zero — idleness is measured by months\_no\_work, and counting it here too would penalise it twice. & (yatra\_months*hours\_week\_yatra + offseason\_months\_worked*hours\_week\_offseason) / (yatra\_months + offseason\_months\_worked) \\
\texttt{work\_income\_pm} & (Yatra income + other work income) divided by 12. & (yatra\_income + non\_yatra\_income)/12 \\
\texttt{emp\_inc\_dep} & 1 if average monthly work income is below six times the national basic food basket. Blank until EMP\_BASKET\_PM is set. & work\_income\_pm < 6*EMP\_BASKET\_PM \\
\texttt{emp\_stab\_dep} & 1 if under 36 months (3 years) in the current job, from the years question. & years\_current\_job < 3 \\
\texttt{emp\_unemp\_dep} & 1 if the respondent looked for work at least once in the last 12 months. Set to 0 where the question was not asked, because it is asked only where there were idle months. & months\_looked\_for\_work > 0 (0 where not asked) \\
\texttt{emp\_sec\_dep} & 1 if not enrolled in any social-security scheme. & social\_security\_any == 0 \\
\texttt{emp\_occ\_dep} & 1 if self-employed without higher education, or wage-employed without a contract (contract code 3 only). & (inlist(employment\_type,1,2) \& higher\_ed==0) | (inlist(employment\_type,3,4) \& contract\_status==3) \\
\texttt{emp\_hours\_dep} & 1 if the average weekly hours across worked months exceed 48. & hours\_week\_annual > 48 \\
\texttt{emp\_intens\_dep} & 1 if at least two of the three intensity demands hold for more than half the working day. & (intens\_speed==1) + (intens\_deadline==1) + (intens\_time==1) >= 2 \\
\texttt{emp\_posture\_dep} & 1 if at least two of the three posture demands hold for more than half the working day. & (posture\_position==1) + (posture\_loads==1) + (posture\_repetitive==1) >= 2 \\
\texttt{emp\_phys\_dep} & 1 if exposed to loud noise OR extreme heat or cold for more than half the working day. & phys\_noise==1 | phys\_temperature==1 \\
\texttt{emp\_score} & Apablaza weights: low earnings 1/4; tenure, unemployment risk, social security and occupational status 1/8 each; the four conditions items 1/16 each. Blank while EMP\_BASKET\_PM is unset. & (1/4)*emp\_inc\_dep + (1/8)*(emp\_stab\_dep+emp\_unemp\_dep+emp\_sec\_dep+emp\_occ\_dep) + (1/16)*(emp\_hours\_dep+emp\_intens\_dep+emp\_posture\_dep+emp\_phys\_dep) \\
\texttt{emp\_poor} & The cut-off of 1/3 is a convention carried over from the MPI, not a figure in the employment table. Check it against Apablaza's own cut-off before reporting. & emp\_score >= 1/3 \\
\texttt{tk\_regular\_n} & Count of task answers equal to 1. & count of tk\_* == 1 \\
\texttt{tk\_prior\_n} & Count of task answers equal to 2. & count of tk\_* == 2 \\
\texttt{other\_act\_1} & 1 if 'Pony/mule worker (works the animals, does not own them)' was ticked in the check-all list of other paid activities. & other\_activity\_types contains 1 \\
\texttt{other\_act\_2} & 1 if 'Pony/mule owner (owns the animals)' was ticked in the check-all list of other paid activities. & other\_activity\_types contains 2 \\
\texttt{other\_act\_3} & 1 if 'Porter (carries on own back/shoulders - goods, luggage, or a person in a kandi)' was ticked in the check-all list of other paid activities. & other\_activity\_types contains 3 \\
\texttt{other\_act\_4} & 1 if 'Palki / dandi bearer (carries a person on a palanquin, as part of a team)' was ticked in the check-all list of other paid activities. & other\_activity\_types contains 4 \\
\texttt{other\_act\_5} & 1 if 'Dhaba, tea stall or food stall WORKER (prepared food or tea)' was ticked in the check-all list of other paid activities. & other\_activity\_types contains 5 \\
\texttt{other\_act\_6} & 1 if 'Dhaba, tea stall or food stall OWNER (prepared food or tea)' was ticked in the check-all list of other paid activities. & other\_activity\_types contains 6 \\
\texttt{other\_act\_7} & 1 if 'Shop WORKER (goods: prasad, puja items, clothes, general store)' was ticked in the check-all list of other paid activities. & other\_activity\_types contains 7 \\
\texttt{other\_act\_8} & 1 if 'Shop OWNER (goods, not prepared food)' was ticked in the check-all list of other paid activities. & other\_activity\_types contains 8 \\
\texttt{other\_act\_9} & 1 if 'Hotel / lodge WORKER' was ticked in the check-all list of other paid activities. & other\_activity\_types contains 9 \\
\texttt{other\_act\_10} & 1 if 'Hotel / lodge OWNER' was ticked in the check-all list of other paid activities. & other\_activity\_types contains 10 \\
\texttt{other\_act\_11} & 1 if 'Driver (any motor vehicle)' was ticked in the check-all list of other paid activities. & other\_activity\_types contains 11 \\
\texttt{other\_act\_12} & 1 if 'Guide' was ticked in the check-all list of other paid activities. & other\_activity\_types contains 12 \\
\texttt{other\_act\_13} & 1 if 'Wage labourer (construction, loading, odd jobs)' was ticked in the check-all list of other paid activities. & other\_activity\_types contains 13 \\
\texttt{other\_act\_14} & 1 if 'Other' was ticked in the check-all list of other paid activities. & other\_activity\_types contains 14 \\
\texttt{best\_alt\_occupation} & Of the 13 occupations, the one (other than the respondent's own) whose task profile is closest to the respondent's. Worker's own task profile uses codes 1 AND 2 (does it now, or has done it before elsewhere) — prior capability outside the current job is exactly what a transfer question is about. Destination profiles are the occupation-level share of workers doing each task (Gathmann and Schonberg's q\_oj). & argmax over j != own occupation of cosine(worker profile, q\_j) \\
\texttt{task\_cover\_best} & Coverage: of the tasks that alternative occupation requires, the share this worker can already perform. Ormiston's t\_ji direction — can the worker step in? & sum(q\_best * worker) / sum(q\_best) \\
\texttt{task\_retain\_best} & Retention: of the tasks this worker can perform, the share that alternative occupation would actually use. Ormiston's t\_ij direction — how much of their skill survives the move? & sum(q\_best * worker) / sum(worker) \\
\texttt{skill\_move\_type} & Direction of the move implied by the coverage/retention pair, both split at 0.5: high-high lateral (move needs little retraining); low coverage but high retention means the destination needs skills the worker lacks while their own still count, i.e. UPSKILLING; high coverage but low retention means the worker has more than the destination uses and much of it goes idle, i.e. DOWNSKILLING; low-low means the skill sets barely overlap, i.e. RESKILLING. Worker's own task profile uses codes 1 AND 2 (does it now, or has done it before elsewhere) — prior capability outside the current job is exactly what a transfer question is about. Destination profiles are the occupation-level share of workers doing each task (Gathmann and Schonberg's q\_oj). & from (task\_cover\_best, task\_retain\_best), each split at 0.5 \\
\bottomrule
\end{longtable}
}

\section*{5. Apablaza et al. (2026) questions: what we asked, changed or dropped}
Module K uses the questions in Appendix 2 of the paper. Their own skip logic is kept: wage workers only for the contract and leave questions; extra hours only if the worker wants more work; the search question only if some month had no work.
The paper asks every household member aged 5 or more. We ask the respondent only (no household roster), so an employment index for the whole household cannot be built.
{\scriptsize
\begin{longtable}{p{5cm}p{4.6cm}p{6.7cm}}
\toprule
\textbf{Apablaza question} & \textbf{Our variable} & \textbf{Note} \\
\midrule
\endfirsthead
\toprule
\textbf{Apablaza question} & \textbf{Our variable} & \textbf{Note} \\
\midrule
\endhead
Q1 household size; Q2 sex, age & \texttt{hhsize, female, age} & Module A \\
Q3 disability & \texttt{not asked} & Respondents are working; not needed for the five domains \\
Q4 birthplace, main earner & \texttt{origin, n\_earners} & Modules A and D \\
Q5 situation last month & \texttt{job\_situation} & Shortened to categories that fit a working respondent \\
Q6 occupation (open) & \texttt{occupation} & Card with 13 groups (the quota variable), not an open answer \\
Q7 activity of the workplace & \texttt{not asked} & Implied by the occupation group \\
Q8 status in the job & \texttt{employment\_type} & PLFS four-way status instead of the paper's seven categories \\
Q9 permanent / seasonal / casual & \texttt{job\_permanence} & Probation and fixed-term merged \\
Q10 years in this job & \texttt{years\_current\_job} & Whole years, so a test for under 6 months is not possible \\
Q11 signed contract & \texttt{contract\_status} & Wage workers only \\
Q12 registered workplace & \texttt{workplace\_registered} & Widened to local registrations (union, Yatra registration, trade licence, GST) \\
Q13 hours in a normal week & \texttt{hours\_day\_yatra, days\_week\_yatra} & Two easier questions, multiplied afterwards \\
Q14 monthly earnings & \texttt{income\_m1 to income\_m12} & Asked month by month (section 6) \\
Q15 earnings range if unknown & \texttt{not asked} & Amounts are asked for every month, so no range is needed \\
Q16 pension & \texttt{pension\_contrib} & As in the paper \\
Q17 health insurance through work & \texttt{work\_health\_ins} & As in the paper \\
Q18 leave rights & \texttt{leave\_rights} & Wage workers only \\
Q19 ever injured; Q20 injury or death at workplace, last 12 months & \texttt{injured\_ever, workplace\_injury\_12m} & The paper's evidence on hazards \\
Q21 wants more work; Q22 how many hours & \texttt{wants\_more\_work, more\_hours\_week} & As in the paper \\
Q23 weeks looking for work & \texttt{looked\_for\_work\_idle} & Tied to the months with no work in the calendar \\
Q24 first job; Q25 why lost the last job & \texttt{not asked} & Apply to the unemployed only; prev\_occ\_reason covers job changes \\
\bottomrule
\end{longtable}
}

The five deprivation indicators (access, compensation, security, stability, working conditions) are built from these raw answers in Stata, using the thresholds in Table 3 of the paper.
Where the paper gives a choice, the fallback is used for pay (67 percent of the median) because the international-dollar version needs a conversion factor we have not verified.

\section*{6. What changed from the earlier draft, and why}
\begin{itemize}[leftmargin=1.2em,itemsep=2pt]
\item \textbf{The 12-month calendar is back.} The worker gives the main activity for each month and the earnings for each month with work. Earnings are asked month by month, so that a single typical figure is not rounded (heaping) and monthly earnings are not guessed from an annual total.
Because it is the original design, the constructed income, the months of work and the seasonality measure now reproduce the earlier pipeline exactly. Months with no work are filled with 0 by the tablet, so the calendar costs about 4.5 minutes.
\item \textbf{Licences and certificates are gone.} People drive, wire and guide without holding a licence, so holding one says little about what they can do. The block asks what people do (the three-answer grid) and adds three factual driving questions, asked only of those who have driven: motorcycle, car, truck or bus, with or without a licence. Entry papers can be studied in the second paper if needed.
\item \textbf{The consumption module now asks a Yatra-season figure AND an off-season figure for every routine item} (9 pairs: staples, perishables, home-grown/gifted food, eating out, fuel, toiletries, transport/communication, rent, non-hospital medical), instead of one ``last 30 days'' figure. Workers here are migrants: fieldwork happens during the Yatra season, so a single recent-recall question would describe only the season the interview happens to fall in, not the whole year. The item groupings and which items get a long (365-day) recall instead still follow HCES 2022-23, IHDS-II and the Nigeria GHS-Panel (read directly; see section 8).
\item \textbf{Remittances are split the same way} (Module C): a usual monthly amount sent and received, for the Yatra season and for the off-season, instead of one annual figure.
\item \textbf{The annual welfare measure is now a genuine annual average}, not an on-season snapshot: each seasonal item is weighted by the number of Yatra-season months and off-season months from the work calendar, exactly as income already was. In the test data this moves \texttt{cons\_pc\_pm} down and the poverty rate up, because the earlier design implicitly treated on-season spending as if it held all year; see the note in \texttt{checks/check\_report.txt}.
\item \textbf{Two new consumption variables for robustness:} a narrow measure (30-day items only, as in Dercon and Krishnan 2000) and a per-adult-equivalent measure (needs the count of children under 15, now asked).
\item \textbf{Four opinion questions on a 1 to 5 scale were removed}, in line with the no-rating rule. Only the three-way ropeway view remains.
\item \textbf{Dropped for length:} remittance frequency and \texttt{prev\_offseason\_work}. The season start, end and length, and the main off-season activity, are now built from the calendar instead of being asked.
\item \textbf{Added:} sex, native language, earners, children under 15 and school-age children, the way the household coped with a shock, direct exposure to the trek, and the Apablaza questions.
\item \textbf{Food security and disability are new} (Modules H and I): a long-term illness/disability flag and an unmet-healthcare-need question (VASyR 2025 / Washington Group short set), plus the standard WFP reduced Coping Strategy Index. This was a genuine gap: a vulnerability-to-poverty instrument with no food-security module. Modelled on the VASyR/Lyons, Kass-Hanna and Montoya Castano (2023) multidimensional livelihood index, read directly (see \texttt{replication/vasyr\_example/}); a full dietary-diversity count was left out to stay inside the 45-minute limit.
\item \textbf{The rCSI is a Yatra-season/off-season pair, like the consumption module}: the first version asked each of the five WFP coping items once, over ``the last 7 days'' -- fieldwork happens during the Yatra season, so that could only describe on-site coping at the work site, not the household's usual pattern (which may look different off-season, with different members present). Each item is now asked as a usual-week day-count for each season, combined into an annual-average rCSI the same way income and consumption already are.
\item \textbf{Household-head and family-structure items are new} (Module A): relationship to the household head and the head's sex (auto-filled with the respondent's own sex when they are the head), and whether the family is nuclear, joint or single-member. \texttt{hhsize} is now asked as household membership directly (``how many members live in your household'') rather than through a shared-kitchen framing.
\item \textbf{A direct residence-migration question is new} (Module D): apart from the seasonal Yatra-work absence, whether the respondent's usual place of residence for most of the year differs from their family's native/home place -- distinct from \texttt{origin} (where they are originally from) and from the NSS short-term-migrant item (a specific 15-day-to-6-month absence threshold). Asked only of non-local-origin respondents.
\item \textbf{The interview site is no longer named.} \texttt{location\_cluster} now has generic labels (Cluster 1-4) instead of the actual place names: enumerators read the site off a printed crosswalk card kept by the research team, not a name on the tablet, so no exported record carries a named interview site. The underlying codes and the health-access-tier logic built from them are unchanged. On the Kobo build (\texttt{questionnaire/build\_xlsform.py}), the enumerator-name question is dropped the same way -- each enumerator has their own Kobo login, which records who submitted automatically -- and GPS is captured silently in the background the moment consent is given, with no on-screen question.
\item \textbf{Tasks block: 12 tasks} (from 27), the ones the provisional destination jobs need most. They fully cover 7 of the 13 destination jobs. The task-transfer paper will need a longer list and employer information.
\item \textbf{Field team is 4 enumerators} in the worked example (the earlier file used 6).
\end{itemize}

\section*{7. Points still open}
\begin{itemize}[leftmargin=1.2em,itemsep=2pt]
\item \textbf{Timing is untested.} The 12-month calendar and the 25-question Apablaza set are the two blocks most likely to run over their estimate. Time them in the pretest.
\item \textbf{The working-conditions indicator is high in the test data} (-- in 100 deprived). It is driven by the social-protection route: almost nobody has a pension and most have no health insurance through work. That may well be true of Yatra workers, but check it against the pilot before treating it as a finding.
The cut-off of 2 of 5 gives -- in 100 employment-poor; a cut-off of 3 gives --.
\item \textbf{Tenure is asked in whole years,} so the paper's under-6-month test for wage workers is not possible; casual jobs supply the variation instead.
\item \textbf{Real quotas.} Occupation quotas must be reset against a real enumeration before fieldwork.
\end{itemize}

\section*{8. The consumption module, read against the three source questionnaires}
Most of the vulnerability-to-poverty papers in the project folder use an existing national survey and report only how consumption was totalled, not the question list itself.
So the three questionnaires were read directly: \textbf{HCES 2022-23} (MoSPI, Appendix A: the source of our poverty line), \textbf{IHDS-II} (the Income and Social Capital Questionnaire, question 14), and the \textbf{Nigeria GHS-Panel} (Wave 3 Household Questionnaire, Sections 10B and 11, the survey behind Eze and Iheonu 2025).
\subsection*{8.1 What each source does}
{\scriptsize
\begin{longtable}{p{3.6cm}p{13.6cm}}
\toprule
\textbf{Source} & \textbf{Recall design} \\
\midrule
\endfirsthead
\toprule
\textbf{Source} & \textbf{Recall design} \\
\midrule
\endhead
HCES 2022-23 (India) & A 'modified mixed reference period': 30 days for cereals, pulses, sugar and salt (Sections 5.1-5.3); 7 days for milk, vegetables, fruit, meat/fish/egg, oil, spices and processed food eaten out (Sections 6-7); 30 days for fuel, toiletries, conveyance, phone and rent (Sections 8-11); 365 days for clothing, footwear, education, hospitalisation and durables (Sections 10, 13-14); 30 days for non-hospitalisation medical (Section 10.3). Insurance is not part of consumption. \\
IHDS-II (India) & A flatter design: 30 days for all 33 food and routine items together (questions 14.1-14.33, including phone, transport and non-hospitalisation medical); 365 days for 19 further items (14.34-14.52: hospitalisation, education, clothing, footwear, durables, and also insurance premiums and vacations, which it counts as consumption). \\
Nigeria GHS-Panel Wave 3 & Four recall tiers: 7 days for food and a short non-food list (tobacco, newspapers, transport); 30 days for fuel, electricity, toiletries, phone/internet, rent; 6 months for clothing, footwear and small household items (shorter than the two India surveys); 12 months for durables, insurance and ceremonies. \\
\bottomrule
\end{longtable}
}

\subsection*{8.2 Where the three sources agree, and where they do not}
\begin{itemize}[leftmargin=1.2em,itemsep=2pt]
\item \textbf{All three agree on 30 days} for fuel and light, toiletries and household consumables, transport and phone charges, and rent.
\item \textbf{All three agree on 12 months} for durable goods (furniture, utensils, appliances).
\item \textbf{HCES and IHDS-II agree exactly} on splitting medical spending into non-hospitalisation (30 days) and hospitalisation (365 days), and on 365 days for clothing, footwear and education. The Nigeria survey instead asks one combined 6-month medical item, and 6 months (not 12) for clothing.
\item \textbf{HCES and IHDS-II disagree on food.} HCES uses 7 days for perishables (milk, vegetables, fruit, meat, oil, spices) to limit recall decay on frequent purchases; IHDS-II asks the same items at 30 days, in one simpler module. The Nigeria survey uses 7 days for all food.
\item \textbf{Only IHDS-II and the Nigeria survey treat insurance premiums as consumption.} HCES does not, which matches standard national-accounts practice (insurance is a financial item). We follow HCES: our insurance question sits in Module G, not in the consumption total.
\item \textbf{HCES and IHDS-II record the source of each food item} (purchased, home-grown, gift) as a column against every single item. That would need a separate question for every food item, so, following Calvo and Dercon (2007) and Eze and Iheonu (2025), home-grown and gifted food is asked as one aggregate question instead.
\end{itemize}
\subsection*{8.3 Our final choice, and why}
We follow \textbf{HCES 2022-23} for which items get a short recall and which get a 365-day recall (staples, fuel, toiletries, transport, rent and non-hospital medical at a routine recall; clothing, education, hospital and durables at 365 days), because HCES is also the source of our poverty line (Sethu et al. 2024). Within the routine group we do not copy HCES's actual-recall design (its 7-day window for perishables, versus 30 days for staples): our workers are migrants, so the bigger source of error is location and season, not recall decay on one item. Every routine item is instead asked as a usual monthly amount, once for the Yatra season and once for the off-season (section 6), which a single actual-recall question cannot give. Every item in Module E of this questionnaire (section 3) carries its exact source citation.
\subsection*{8.4 Other papers in the folder}
{\scriptsize
\begin{longtable}{p{3.6cm}p{7.6cm}p{5.1cm}}
\toprule
\textbf{Source} & \textbf{Consumption measure} & \textbf{Use in our list} \\
\midrule
\endfirsthead
\toprule
\textbf{Source} & \textbf{Consumption measure} & \textbf{Use in our list} \\
\midrule
\endhead
Calvo and Dercon 2007 (Ethiopia) & Total value of food and non-food consumption, from purchases, own harvest and gifts. & Source for asking home-grown and gifted food as one aggregate question. \\
Dercon and Krishnan 2000 (Ethiopia) & Monthly consumption: food plus a deliberately narrow non-food list, to avoid seasonal swings from large expenditures. & Source of the narrow (30-day only) consumption measure kept as a robustness check. \\
Azeem et al. 2016 (Punjab, MICS) & Quantities of foods consumed and monetary non-food spending, valued at weekly prices. & We ask values, not quantities, to keep the interview short. \\
Azeem et al. 2018 (Pakistan); Awuni et al. 2023 (Ghana) & Food intake converted to an adult-equivalent basis (Claro et al. 2010 scale). & Source of the adult-equivalent consumption measure kept as a robustness check; needs the count of children under 15. \\
Lyons et al. 2023 (VASyR, Lebanon) & A short expenditure module (food total, rent, water, income) run at national scale. & Shows a short module of this kind can work in the field. \\
Khosla and Jena 2022; Jha et al. 2018 (rural India, using IHDS/NSS) & Monthly per capita consumption expenditure; item list not given in the paper (the source survey was read directly instead). & Confirms the outcome variable is per-capita monthly consumption. \\
Ding 2022; Wei et al. 2023; Ceesay and Morelli 2026 & Income-based, or income and consumption used interchangeably, in small samples. & Not comparable; we keep consumption, not income, as the welfare measure. \\
Zheng and Ma 2023 (China, CLES) & Ten food groups counted over the reference week, for dietary diversity, not value. & Not used; our aim is a value-based poverty line, not a diet-diversity score. \\
\bottomrule
\end{longtable}
}

\section*{9. Code lists}
{\scriptsize
\begin{longtable}{p{2cm}p{14cm}}
\toprule
\textbf{List} & \textbf{Codes} \\
\midrule
\endfirsthead
\toprule
\textbf{List} & \textbf{Codes} \\
\midrule
\endhead
\texttt{yn} & 0 No; 1 Yes \\
\texttt{ropexp} & 1 More work for me; 2 Less work for me; 3 About the same; 97 Don't know \\
\texttt{ropjobs} & 1 People from these villages; 2 People from outside the area; 3 Nobody will get jobs; 97 Don't know \\
\texttt{obsenv} & 1 Normal crowd, normal weather; 2 High crowd or peak hours; 3 Rain or extreme weather; 4 Both high crowd and bad weather \\
\texttt{permdk} & 1 Yes, it is allowed; 0 No, it is not allowed; 2 Allowed only with a permit or a fee; 97 Don't know \\
\texttt{cpruse} & 0 No; 1 Yes; 2 No, it is not allowed here; 97 Don't know \\
\texttt{assets} & 1 Television; 2 Radio; 3 Bicycle; 4 Motorcycle or scooter; 5 Car, jeep or truck; 6 Telephone of any kind (mobile or landline); 7 Computer or laptop; 8 Cart pulled by an animal; 9 Refrigerator \\
\texttt{livestock} & 1 Cows or buffaloes; 2 Goats or sheep; 3 Ponies or mules \\
\texttt{offact} & 2 Farming (crops); 3 Livestock (animals); 4 Casual or daily wage labour; 5 Construction work; 6 Own small shop, stall or trade; 7 Salaried job; 96 Some other work; 0 There were no such months \\
\texttt{gpserr} & 1 Location not allowed in the browser; 2 No fix obtained; 3 Coarse fix only, over 200 m -- network or IP location, not GPS \\
\texttt{wageemp} & 1 Private company; 2 Government or public body; 3 A household (domestic work); 4 A contractor or agent; 5 An individual employer; 97 Don't know \\
\texttt{earnrange} & 1 Under Rs 5,000; 2 Rs 5,000 to 10,000; 3 Rs 10,000 to 20,000; 4 Rs 20,000 or more; 97 Don't know \\
\texttt{cprwater} & 1 Yes, free to use; 2 A source exists but we must pay; 3 No source nearby; 97 Don't know \\
\texttt{cprgov} & 1 Village head or panchayat; 2 Forest department; 3 Local committee or community group; 4 Nobody decides; it is open to all; 97 Don't know \\
\texttt{sex} & 0 Male; 1 Female \\
\texttt{enum} & 1 Raman; 2 Rishit; 3 Tanmay; 4 Anuj \\
\texttt{nativelang} & 1 Garhwali; 2 Kumaoni; 3 Hindi; 4 Nepali; 5 Bhojpuri; 6 Assamese; 7 Bengali; 8 Bodo; 9 Dogri; 10 Gujarati; 11 Kannada; 12 Kashmiri; 13 Konkani; 14 Maithili; 15 Malayalam; 16 Manipuri (Meitei); 17 Marathi; 18 Odia; 19 Punjabi; 20 Sanskrit; 21 Santali; 22 Sindhi; 23 Tamil; 24 Telugu; 25 Urdu; 96 Other Indian language; 97 Other / foreign language \\
\texttt{state} & 1 Andhra Pradesh; 2 Arunachal Pradesh; 3 Assam; 4 Bihar; 5 Chhattisgarh; 6 Goa; 7 Gujarat; 8 Haryana; 9 Himachal Pradesh; 10 Jharkhand; 11 Karnataka; 12 Kerala; 13 Madhya Pradesh; 14 Maharashtra; 15 Manipur; 16 Meghalaya; 17 Mizoram; 18 Nagaland; 19 Odisha; 20 Punjab; 21 Rajasthan; 22 Sikkim; 23 Tamil Nadu; 24 Telangana; 25 Tripura; 26 Uttar Pradesh; 28 West Bengal; 29 Andaman and Nicobar Islands; 30 Chandigarh; 31 Dadra and Nagar Haveli and Daman and Diu; 32 Delhi; 33 Jammu and Kashmir; 34 Ladakh; 35 Lakshadweep; 36 Puducherry \\
\texttt{ruralurban} & 1 Village (rural); 2 Town or city (urban) \\
\texttt{landunit} & 1 Nali; 2 Bigha; 3 Acres; 4 Hectares; 5 No land \\
\texttt{loanpurpose} & 1 Medical or hospital costs; 2 Buying an animal, vehicle or equipment for work; 3 Shop stock or business; 4 Food or daily household needs; 5 A wedding, funeral or ceremony; 6 House building or repair; 7 Education; 8 Repaying another loan; 9 Something else \\
\texttt{collat} & 0 Nothing was given as security; 1 Jewellery or ornaments; 2 Land; 3 Animals; 4 A vehicle; 5 Shop stock or business goods; 6 House or building; 7 Something else \\
\texttt{recall} & 1 Month by month; 0 One total for the whole year \\
\texttt{hohrel} & 1 Self (respondent is the head); 2 Husband; 3 Wife; 4 Father; 5 Mother; 6 Son; 7 Daughter; 8 Brother; 9 Sister; 10 Other male relative; 11 Other female relative; 12 Other, non-relative (male); 13 Other, non-relative (female) \\
\texttt{famstruct} & 1 Nuclear (you, spouse and unmarried children only); 2 Joint (living with parents, married siblings or other extended family); 3 Single-member household \\
\texttt{marital} & 1 Currently married; 2 Never married; 3 Widowed/divorced/separated \\
\texttt{edu2} & 1 No formal education; 2 Some schooling, did not complete primary; 3 Completed primary but not secondary; 4 Completed secondary or higher; 98 Prefer not to say \\
\texttt{occ} & 1 Pony/mule worker (works the animals, does not own them); 2 Pony/mule owner (owns the animals); 3 Porter (carries on own back/shoulders - goods, luggage, or a person in a kandi); 4 Palki / dandi bearer (carries a person on a palanquin, as part of a team); 5 Dhaba, tea stall or food stall WORKER (prepared food or tea); 6 Dhaba, tea stall or food stall OWNER (prepared food or tea); 7 Shop WORKER (goods: prasad, puja items, clothes, general store); 8 Shop OWNER (goods, not prepared food); 9 Hotel / lodge WORKER; 10 Hotel / lodge OWNER; 11 Driver (any motor vehicle); 12 Guide; 13 Wage labourer (construction, loading, odd jobs); 14 Other \\
\texttt{emptype} & 1 Own-account (self-employed, no hired workers); 2 Employer/business owner (self-employed); 3 Regular wage or salary; 4 Casual or daily wage labour; 5 Unpaid family worker \\
\texttt{site} & 1 Kedarnath (Gaurikund to Kedarnath); 2 Hemkund Sahib (Govindghat to Hemkund) \\
\texttt{accomhere} & 1 At my own usual home (I live here); 2 Rented room or house; 3 Room or dormitory the employer provides; 4 Sleeps at the shop, dhaba or workplace; 5 Tent or temporary shelter; 6 Lodge, dharamshala or ashram; 7 In the open, or a verandah; 8 Somewhere else \\
\texttt{wgdiff} & 1 No difficulty; 2 Some difficulty; 3 A lot of difficulty; 4 Cannot do it at all \\
\texttt{portable} & 1 Yes, can draw it here; 2 No, only at the home place; 3 Has never tried; 4 We have no ration card; 97 Does not know \\
\texttt{month} & 1 January; 2 February; 3 March; 4 April; 5 May; 6 June; 7 July; 8 August; 9 September; 10 October; 11 November; 12 December \\
\texttt{prevreason} & 1 Better income; 2 Lost the previous work; 3 Work ended with the season; 4 Family reasons; 5 Health or injury; 6 Other \\
\texttt{origin} & 1 Local (same district); 2 Other Uttarakhand district; 3 Other Indian state; 4 Nepal \\
\texttt{referral} & 1 A family member already working here; 2 A friend or someone from the village; 3 A thekedar or contractor I now work for; 4 An agent or middleman who placed me (usually for a fee); 5 No one; I found it myself; 6 The employer called me directly; 7 Someone else \\
\texttt{floor} & 1 Mud/kaccha; 2 Cement/mud-cement; 3 Tile/mosaic/marble \\
\texttt{roof} & 1 Thatch/wood/mud; 2 Tin/GI sheet; 3 Concrete/RCC \\
\texttt{wall} & 1 Mud, thatch, bamboo or other natural material; 2 Unburnt brick, wood or tin; 3 Burnt brick, cement, concrete or stone \\
\texttt{toilet} & 1 No toilet / open defecation; 2 Pit latrine without slab or open pit; 3 Improved toilet, but shared with other households; 4 Improved toilet, used only by this household \\
\texttt{water} & 1 Piped into the house or yard; 2 Public tap or standpipe; 3 Handpump, tubewell or borewell; 4 Protected well or protected spring; 5 Rainwater collection; 6 Bottled, packaged or community RO; 7 Unprotected well or unprotected spring; 8 River, stream, pond or canal; 9 Tanker truck or cart with drum \\
\texttt{fuel} & 1 LPG or cylinder gas; 2 Piped natural gas; 3 Electricity; 4 Biogas (gobar gas plant); 5 Kerosene; 6 Firewood; 7 Dung cakes (gobar); 8 Crop residue, straw or shrubs; 9 Charcoal; 10 Coal or lignite \\
\texttt{comereason} & 1 No work at home; 2 Pay is better here; 3 Family or people from my village were already here; 4 Land at home is too little to live on; 5 Debt to repay; 6 A contractor or agent brought me; 7 Married into / moved with family; 8 Some other reason \\
\texttt{govtscheme} & 1 Ration card / PDS food grain; 2 MGNREGA work or job card; 3 Old-age, widow or disability pension; 4 PM-KISAN (farmer cash transfer); 5 Ujjwala LPG connection; 6 PM-SYM or Atal Pension Yojana; 7 Ayushman Bharat / health insurance card; 8 Housing scheme (PMAY or state); 9 Some other scheme; 10 None of these \\
\texttt{remitmode} & 1 From my own phone (UPI or net banking); 2 In person at a bank branch or bank mitra / CSP; 3 Money order or post office; 4 Sent with someone going home; 5 Carried it myself; 6 Through an agent or middleman; 7 Other \\
\texttt{credit} & 1 Nationalised or public-sector bank; 2 Private bank; 3 Cooperative bank or RRB; 4 Microfinance institution or SHG; 5 Moneylender; 6 Relative or friend; 7 Employer or contractor (advance); 8 Other \\
\texttt{distress} & 9 Nothing of this kind happened; 1 Serious illness or death of an earning member; 2 Crop failure or livestock loss; 3 Damage from a natural disaster (flood, landslide, fire); 4 Loss of business, goods or assets; 5 Road or bridge blocked, route closed; 6 Yatra stopped early or badly disrupted; 7 Sharp fall in customers or prices; 8 Lost the work or the work stopped \\
\texttt{coping} & 1 Used savings; 2 Borrowed money; 3 Sold or pawned assets; 4 Cut consumption; 5 Help from relatives/community; 6 Paid it out of normal earnings, nothing given up; 7 Something else (say what) \\
\texttt{birthattend} & 1 Doctor; 2 Nurse, ANM or LHV; 3 ASHA or Anganwadi worker only; 4 Dai (traditional birth attendant); 5 Nobody trained; family only \\
\texttt{stance} & 1 Support; 2 Neutral; 3 Oppose; 98 Prefer not to say \\
\texttt{task3} & 1 Yes, regularly in my main Yatra work; 2 Not in my main work, but done before elsewhere; 3 Never done \\
\texttt{tier} & 1 Poor; 2 Moderate; 3 Good \\
\texttt{activity} & 1 Yatra work (working the season); 9 Getting ready for the season (not earning from it yet); 2 Farming (crops); 3 Livestock (animals); 4 Casual or daily wage labour; 5 Construction work; 6 Own small shop, stall or trade; 7 Salaried job; 8 No paid work \\
\texttt{yndk} & 0 No; 1 Yes; 97 Don't know \\
\texttt{jobsit} & 1 Works for pay full-time; 2 Works for pay part-time or occasional jobs; 3 Studies and works; 4 Only studies; 5 Being trained for work only; 6 Retired or pensioned; 7 Unpaid household tasks or caring for others; 8 Unemployed, actively seeking work; 9 Sick or disabled, cannot work; 10 Neither studying, working nor seeking work; 97 Does not know \\
\texttt{jobperm} & 1 Permanent; 2 Seasonal or temporary; 3 Occasional or casual; 4 On probation; 5 Fixed-term \\
\texttt{contract} & 1 Yes, signed; 2 Yes, but not yet signed; 3 No contract \\
\texttt{pension} & 1 Yes, employer deducts it; 2 Yes, contributes voluntarily; 3 No \\
\texttt{workins} & 1 Yes, through work; 2 Only private or other insurance; 3 None; 4 Don't know \\
\texttt{movetype} & 1 Lateral (little retraining needed); 2 Upskilling (destination needs skills the worker lacks); 3 Reskilling (skill sets barely overlap); 4 Downskilling (worker's skills would go unused) \\
\bottomrule
\end{longtable}
}
\end{document}
